\chapter{Quiet and informative markets}\label{ch:coexistence}

\begin{chapterplan}
The question of this chapter: against one and the same incumbent, can the price stay silent in one equilibrium and
recruit a high-cost challenger in another, and where along a path of incumbents does this happen?

After this chapter you can:
\begin{itemize}
\item explain in words why a quiet market and an informative market can each confirm itself against the same
incumbent;
\item derive the condition for a quiet market and the four conditions for coexistence, one step at a time;
\item derive the purchase probability of class \(H\) and entry in the informative market, and show that entry rises
with the spread while class \(H\) mixes;
\item prove the three parts of Proposition~3 with the lemmas on the high class, the low class and the product
structure;
\item read Corollary~2 and Figure~\ref{fig:coexistence-path}, name the marked strengths in their order, and say which
conclusions depend on the tie rules.
\end{itemize}
\end{chapterplan}

% =====================================================================
\section{Two ways for one market to work}\label{sec:coexistence-words}

Chapter~\ref{ch:ccc} compared three economies and found exactly one equilibrium in each. There the strict
conditions (H1) to (H3) fixed the investor's orders before anything else: against the weak incumbent no order could
pay, and against the strong incumbent every correct order paid. This chapter looks at the strengths between those
two. There the trading bound \eqref{eq:trading-bound} does not decide the orders (Section~\ref{sec:trading-bound}),
and the answer depends on what the market expects.

Fix one incumbent. Suppose first that nobody expects the price to recruit a high-cost challenger. Then entry is the
floor \(\rho\) after every order. An informed order earns entry times the spread times the doubt that remains about
the class (Section~\ref{sec:trading-returns}), so with entry at the floor it earns little. If this is below the
trading cost \(k\), the investor stays out. With no informed orders, every order comes from a liquidity trader, the
price says nothing, and the belief stays at the prior. If the prior lies below the preparation threshold, no
high-cost challenger prepares. So nobody has a reason to expect recruitment, and the expectation confirms itself.
This is a \emph{quiet market}.

Suppose instead that everybody expects a purchase to recruit. Then entry after a purchase is one. At a fixed
posterior, this multiplies the return of a purchase by \(1/\rho\). If the larger return covers \(k\), class \(H\)
buys. Its purchases make a purchase good news about class \(H\). If the posterior after a purchase reaches the
threshold \(\tau\), a high-cost challenger does prepare after a purchase. Again the expectation confirms itself.
This is an \emph{informative market}.

Both stories use the same incumbent, the same costs and the same order slot. Only the expectation differs. Each
story is an equilibrium in the sense of Section~\ref{sec:toolkit-games}: given what the others do, no player gains
by a change of its own plan alone. That section named this case as its example of several equilibria.

\begin{keyidea}
Recruitment raises the return of a purchase, and informed purchases are what make recruitment happen. So ``no
trade and no recruitment'' and ``trade and recruitment'' can both be consistent against one incumbent.
\end{keyidea}

\begin{analogy}[a scout and a club]
This analogy is the guide's own. A scout knows whether a young player is good. A tip to a club costs the scout a
fixed fee to send. The scout earns a share of the gain only when the club invites the player to a trial, and a trial
is expensive for the club.

Suppose that clubs ignore tips. Then a tip rarely leads to a trial and earns little, so the scout does not pay the
fee. No tips arrive, and the clubs are right to ignore the tip line. Suppose instead that clubs invite a player after
a tip. Then a tip earns much more, and the scout sends tips about good players. A tip becomes good news, and the
clubs are right to invite.

Where the analogy breaks:
\begin{itemize}
\item The investor is not paid by the challenger. It earns from trade, and market makers put the expected entry
into the price. So the investor earns only on the doubt that remains after its order.
\item A liquidity trader can place the same order. A purchase is never proof of class \(H\).
\item In the model some challengers prepare whatever the price says: the zero-cost challenger, with probability
\(\rho\). So even a silent price leaves the investor a small return. A very wide spread makes this small return
cover \(k\), and then the quiet market ends.
\item When a purchase earns exactly \(k\), the investor may buy with any probability. The model uses this: in an
informative market class \(H\) may mix.
\end{itemize}
\end{analogy}

Lemma~A.4 (Section~\ref{sec:trading-twonumbers}) says that orders, posteriors and entry depend on the auction only
through two numbers: the spread \(\DT\) and the threshold \(\tau\). So the incumbents at which both markets exist form
a region in the plane of these two numbers. Section~4.3 of the paper describes the region in words:
\begin{itemize}
\item the threshold lies above the prior, so an uninformative price does not recruit;
\item the threshold lies within reach of a purchase;
\item the spread is wide enough that a purchase that just reaches the threshold, and recruits, covers \(k\);
\item the spread is narrow enough that no order against an uninformative price earns more than \(k\).
\end{itemize}
Sections~\ref{sec:coexistence-quiet} and~\ref{sec:coexistence-informative} derive each condition.

% =====================================================================
\section{The set-up}\label{sec:coexistence-setup}

This section recalls the symbols of the chapter and adds a few short names from Appendix~A.4 of the paper.
\begin{itemize}
\item The challenger's class is \(j\in\{H,L\}\), with prior \(\pi=\Pr(j=H)\in(0,1)\).
\item With probability \(\lambda\in(0,1)\) the informed investor fills the order slot. It knows the class and places
\(q\in\{-1,0,1\}\). A nonzero order costs \(k>0\). Otherwise a liquidity trader sells, abstains or buys with
probabilities \(\beta\), \(1-2\beta\) and \(\beta\). Appendix~A.4 fixes \(0<\beta<1/2\).
\item A liquidity order in one direction has probability \(b=(1-\lambda)\beta\), a liquidity abstention has
probability \(b_0=(1-\lambda)(1-2\beta)\), and \(u=\lambda/b\).
\item The posterior after the order \(x\) is \(\mu(x)\). The top and bottom posteriors are \(M=M(\pi)\) and
\(m=m(\pi)\) of \eqref{eq:beliefs-bounds}.
\item The preparation cost is zero with probability \(\rho\in(0,1)\) and \(c_H>0\) otherwise. Entry after the order
\(x\) is \(e(x)=\rho+(1-\rho)\,\zeta(P(x))\), so \(e(x)\ge\rho\): the entry floor.
\item The incumbent law \(F\) is continuous, with a positive spread \(\DT=t_H-t_L>0\) and \(g_H>g_L\). The
preparation threshold is \(\tau=(c_H-g_L)/(g_H-g_L)\), equation \eqref{eq:beliefs-threshold}.
\item The gross returns of the two correct orders are \(K_H=e(1)\DT[1-\mu(1)]\) and \(K_L=e(-1)\DT\mu(-1)\),
equation \eqref{eq:trading-returns}. A wrong-signed order earns at most \(-k\) (Lemma~A.3(a)). Class \(H\) buys if
\(K_H>k\) and abstains if \(K_H<k\); at \(K_H=k\) both orders are best (Lemma~A.3(b)). Class \(L\) is the same with
\(K_L\) and a sale.
\item The default tie rules (Section~\ref{sec:model-ties}): an indifferent investor may mix over its best orders,
and an indifferent high-cost challenger prepares.
\end{itemize}

With \(\DT>0\), \(g_H>g_L\) and \(\rho>0\), the price reveals the posterior (Section~\ref{sec:beliefs-reveal}). Under
the default challenger rule, entry follows the posterior (Lemma~A.2(d)):
\[
e(x)=\rho\quad\text{if }\mu(x)<\tau,\qquad e(x)=1\quad\text{if }\mu(x)\ge\tau .
\]

The paper uses two words for equilibria. A \emph{quiet market} is an equilibrium in which both classes abstain,
\(\sigma_H=\sigma_L=\delta_0\). An \emph{informative market} is an equilibrium with informed trade: some class places
a nonzero order with positive probability.

Appendix~A.4 adds three short names:
\[
J=\DT\,(1-\tau),\qquad \DT^{\rm F}=\frac{k}{1-M},\qquad \DT^{\rm N}=\frac{k}{\rho\max\{\pi,1-\pi\}} .
\]
\begin{itemize}
\item \(J\) is the gross return of a purchase whose posterior just reaches the threshold and that recruits: put
\(e(1)=1\) and \(\mu(1)=\tau\) into \(K_H\) (Section~\ref{sec:trading-path}).
\item \(\DT^{\rm F}\) is the spread at which a full purchase that recruits earns exactly \(k\):
\(\DT^{\rm F}(1-M)=k\).
\item \(\DT^{\rm N}\) is the spread at which the better of the two unit orders against an uninformative price, with
entry \(\rho\), earns exactly \(k\). Section~\ref{sec:coexistence-quiet} derives it.
\end{itemize}
The letters follow the paper's labels: F for a full purchase, N for a quiet market.

\subsection{The posterior after a purchase and after a sale}

By Lemma~A.3(a), no equilibrium puts weight on a wrong-signed order. So class \(L\) never buys and class \(H\) never
sells. Write
\[
s_+=\lambda\,\sigma_H(1),\qquad s_-=\lambda\,\sigma_L(-1)
\]
for the probability that the investor is present and buys in class \(H\), and that it is present and sells in
class \(L\). The guide calls them the purchase rate and the sale rate. Both lie in \([0,\lambda]\).

\begin{derivation}[order probabilities]
Equation~\eqref{eq:model-orderprob} gives \(f_j(x)=(1-\lambda)\eta(x)+\lambda\sigma_j(x)\). In class \(H\),
\(\sigma_H(-1)=0\) and \(\lambda\sigma_H(0)=\lambda-s_+\). In class \(L\), \(\sigma_L(1)=0\) and
\(\lambda\sigma_L(0)=\lambda-s_-\). So, on the orders \((-1,0,1)\),
\begin{align*}
f_H &= \bigl(b,\; b_0+\lambda-s_+,\; b+s_+\bigr), \why{liquidity part plus informed part}\\
f_L &= \bigl(b+s_-,\; b_0+\lambda-s_-,\; b\bigr). \why{the same in class \(L\)}
\end{align*}
\end{derivation}

Bayes' rule (Section~\ref{sec:toolkit-bayes}) turns these into posteriors. Write \(\mu_+(s_+)\) for the posterior
after a purchase and \(\mu_-(s_-)\) for the posterior after a sale.

\begin{derivation}[the posteriors]
\begin{align*}
\mu_+(s_+) &= \frac{\pi(b+s_+)}{\pi(b+s_+)+(1-\pi)\,b} \why{Bayes' rule at \(x=1\)}\\
&= \frac{\pi(b+s_+)}{b+\pi s_+} \why{\(\pi b+(1-\pi)b=b\)}\\[3pt]
1-\mu_+(s_+) &= \frac{b+\pi s_+-\pi b-\pi s_+}{b+\pi s_+} \why{common denominator}\\
&= \frac{(1-\pi)\,b}{b+\pi s_+} \why{collect terms}\\[3pt]
\mu_-(s_-) &= \frac{\pi b}{\pi b+(1-\pi)(b+s_-)} \why{Bayes' rule at \(x=-1\)}\\
&= \frac{\pi b}{b+(1-\pi)s_-} \why{\(\pi b+(1-\pi)b=b\)}
\end{align*}
The odds form (Section~\ref{sec:toolkit-odds}) gives the zero order:
\[
\frac{\mu_0}{1-\mu_0}=\frac{\pi}{1-\pi}\cdot\frac{b_0+\lambda-s_+}{b_0+\lambda-s_-}.
\]
\end{derivation}

Each posterior moves in one direction with its rate.
\begin{itemize}
\item In \(1-\mu_+(s)=(1-\pi)b/(b+\pi s)\), the numerator is fixed and the denominator rises with \(s\). So \(\mu_+\)
rises strictly. At \(s=0\) it is \(\pi\). As \(s\) grows without bound, the formula approaches one.
\item In \(\mu_-(s)=\pi b/[b+(1-\pi)s]\), the numerator is fixed and the denominator rises with \(s\). So \(\mu_-\)
falls strictly from \(\pi\) at \(s=0\).
\end{itemize}
At the full rate \(s=\lambda\), the two formulas give the bounds of \eqref{eq:beliefs-bounds}:
\begin{align*}
1-M &= \frac{1-\pi}{\pi(1+u)+1-\pi} \why{one minus \(M(\pi)\)}\\
&= \frac{(1-\pi)\,b}{\pi(b+\lambda)+(1-\pi)\,b} \why{multiply above and below by \(b\); \(bu=\lambda\)}\\
&= \frac{(1-\pi)\,b}{b+\pi\lambda}=1-\mu_+(\lambda) \why{\(\pi b+(1-\pi)b=b\)}
\end{align*}
In the same way, \(m=\pi b/[b+(1-\pi)\lambda]=\mu_-(\lambda)\). So on \([0,\lambda]\) the posterior after a purchase
rises from \(\pi\) to \(M\), and the posterior after a sale falls from \(\pi\) to \(m\).

\begin{intuition}[a purchase is better news when class \(H\) buys more often]
The liquidity purchases are a fixed cover in both classes. Each extra purchase by class \(H\) tilts the pool of
purchases toward class \(H\). So the more often class \(H\) buys, the more a purchase reveals, and the less doubt it
leaves.
\end{intuition}

\subsection{The purchase rate that just reaches the threshold}

Let \(\pi<\tau<1\). At which purchase rate does the posterior after a purchase equal \(\tau\)?

\begin{derivation}[the rate \(s^*\)]
\begin{align*}
\frac{\pi(b+s)}{b+\pi s} &= \tau \why{the posterior equals the threshold}\\
\pi(b+s) &= \tau(b+\pi s) \why{multiply by \(b+\pi s>0\)}\\
\pi b+\pi s &= \tau b+\tau\pi s \why{expand}\\
\pi s-\tau\pi s &= \tau b-\pi b \why{collect the terms in \(s\) on the left}\\
\pi s\,(1-\tau) &= b\,(\tau-\pi) \why{factor}\\
s^* &= \frac{b\,(\tau-\pi)}{\pi\,(1-\tau)} \why{divide by \(\pi(1-\tau)>0\)}
\end{align*}
\end{derivation}

The rate \(s^*\) is positive because \(\tau>\pi\). Since \(\mu_+\) rises strictly:
\begin{itemize}
\item \(\mu_+(s_+)<\tau\) exactly when \(s_+<s^*\);
\item \(s^*\le\lambda\) exactly when \(\tau=\mu_+(s^*)\le\mu_+(\lambda)=M\), and \(s^*<\lambda\) exactly when \(\tau<M\);
\item if \(\tau\ge M\), every \(s_+<\lambda\) has \(\mu_+(s_+)<M\le\tau\), so no purchase at a rate below \(\lambda\)
reaches the threshold.
\end{itemize}
In words, \(s^*\) is the least purchase rate at which a purchase is good enough news to recruit.

% =====================================================================
\section{The quiet market}\label{sec:coexistence-quiet}

In a quiet market both classes abstain, so \(s_+=s_-=0\). This section finds when it is an equilibrium, in the case
\(\pi<\tau\). The proof of Proposition~3(a) treats the case \(\pi\ge\tau\) (Section~\ref{sec:coexistence-proof-a}).

The intuition is short. In a quiet market the price is silent, so the only entry that an informed order can use is
the floor: the zero-cost challenger prepares whatever the price says. A deviation by the investor therefore earns
\(\rho\DT\) times the doubt at the prior. The quiet market survives as long as this floor return stays at or below
\(k\) in both classes.

\begin{derivation}[the quiet market, when \(\pi<\tau\)]
\emph{Step 1: every posterior is the prior.} With \(s_+=s_-=0\), the order probabilities of
Section~\ref{sec:coexistence-setup} agree in the two classes: \(f_H=f_L=(b,\,b_0+\lambda,\,b)\). The odds form gives
\begin{align*}
\frac{\mu(x)}{1-\mu(x)} &= \frac{\pi}{1-\pi}\cdot\frac{f_H(x)}{f_L(x)} \why{odds form of Bayes' rule}\\
&= \frac{\pi}{1-\pi} \why{\(f_H(x)=f_L(x)\)}
\end{align*}
Equal odds give equal probabilities, so \(\mu(x)=\pi\) after every order. This is Step~1 of Part~(i) in
Section~\ref{sec:ccc-proof}.

\emph{Step 2: one price, and entry at the floor.} Every belief is \(\pi<\tau\), so a high-cost challenger stays out,
with no tie, and \(e(x)=\rho\) after every order. The price formula \eqref{eq:beliefs-price} gives the same price
\(t_0+\rho\,[t_L-t_0+\DT\pi]\) after every order. This single price pools all orders, and its belief is \(\pi\).

\emph{Step 3: the deviations.} A deviation by the investor faces the candidate's price and entry
(Section~\ref{sec:model-ties}). After a purchase, the price is the pooled price, with belief \(\pi\) and entry
\(\rho\). So
\begin{align*}
K_H &= e(1)\,\DT\,[1-\mu(1)] \why{equation~\eqref{eq:trading-returns}}\\
&= \rho\,\DT\,(1-\pi) \why{\(e(1)=\rho\), \(\mu(1)=\pi\)}\\[3pt]
K_L &= e(-1)\,\DT\,\mu(-1) \why{equation~\eqref{eq:trading-returns}}\\
&= \rho\,\DT\,\pi \why{\(e(-1)=\rho\), \(\mu(-1)=\pi\)}
\end{align*}

\emph{Step 4: abstention must be a best order.} A wrong-signed order earns at most \(-k<0\). Abstention earns zero.
By Lemma~A.3(b), abstention is a best order of class \(H\) exactly when \(K_H\le k\). If \(K_H<k\), it is the only
best order. If \(K_H=k\), the purchase and abstention are both best, and the default investor rule allows any mix,
so abstention for sure is allowed. The same holds for class \(L\) with \(K_L\).

\emph{Step 5: combine.}
\begin{align*}
&\rho\,\DT\,(1-\pi)\le k\ \text{ and }\ \rho\,\DT\,\pi\le k \why{Steps 3 and 4}\\
&\iff \rho\,\DT\,\max\{\pi,1-\pi\}\le k \why{both hold exactly when the larger one holds}\\
&\iff \DT\le\frac{k}{\rho\max\{\pi,1-\pi\}}=\DT^{\rm N} \why{divide by \(\rho\max\{\pi,1-\pi\}>0\)}
\end{align*}
\end{derivation}

So, when \(\pi<\tau\),
\begin{equation}\label{eq:coexistence-quiet}
\text{a quiet market is an equilibrium}\iff\DT\le\DT^{\rm N}=\frac{k}{\rho\max\{\pi,1-\pi\}} .
\end{equation}
There is then exactly one quiet market: its entry is \(\rho\) after every order, and the price formula fixes its
price.

\paragraph{What this says.} The quiet market needs a narrow spread. It does not need a wide one, because nothing in it
uses the spread to pay for trade. Test the extremes.
\begin{itemize}
\item As \(\rho\) falls toward zero, \(\DT^{\rm N}\) grows without bound. Without the zero-cost challenger, a silent
price would give an informed order nothing, and a quiet market would survive every spread.
\item At the prior \(1/2\), \(\max\{\pi,1-\pi\}=1/2\) and \(\DT^{\rm N}=2k/\rho\). This is the largest value of
\(\DT^{\rm N}\) over all priors, because \(\max\{\pi,1-\pi\}\ge1/2\).
\item For a prior near zero or one, \(\max\{\pi,1-\pi\}\) is near one and \(\DT^{\rm N}\) is near \(k/\rho\).
\end{itemize}

\begin{watchout}
In Step~3 a purchase does not recruit, because the price after a purchase is the pooled price with belief \(\pi\). A
deviation faces the candidate's prices and entry. It does not bring the informative market with it. The comparison of
the two markets is a comparison of two equilibria, each solved on its own.
\end{watchout}

\begin{checkyourself}\label{q:coexistence-binds}
Let \(\pi>1/2\) and \(\pi<\tau\). In a quiet market, which of \(K_H\) and \(K_L\) is larger, and which class is
closer to trade? In the trading bound of Section~\ref{sec:trading-bound}, the class that the prior favours binds.
Explain in words why the quiet market gives the opposite answer.
\end{checkyourself}

% =====================================================================
\section{The informative market}\label{sec:coexistence-informative}

Keep \(\pi<\tau\), and let the quiet market be an equilibrium: \(\DT\le\DT^{\rm N}\). This section finds every
informative market that can exist next to it. The argument has four steps: the low class, the zero order, the
purchase, and the purchase rate.

\paragraph{Step 1: class \(L\) abstains.} A sale is bad news for class \(H\), so it never recruits. It earns at floor
entry, and the quiet-market condition says that orders at floor entry do not pay, even at the prior. A sale that the
market expects pays even less, because it leaves less doubt.

\begin{derivation}[class \(L\) abstains]
The posterior after a sale is \(\mu_-(s_-)\le\mu_-(0)=\pi<\tau\), so \(e(-1)=\rho\). Suppose \(s_->0\). Then
\begin{align*}
K_L &= e(-1)\,\DT\,\mu_-(s_-) \why{equation~\eqref{eq:trading-returns}}\\
&= \rho\,\DT\,\mu_-(s_-) \why{\(e(-1)=\rho\)}\\
&< \rho\,\DT\,\pi \why{\(\mu_-(s_-)<\pi\) for \(s_->0\); \(\rho\DT>0\)}\\
&\le \rho\,\DT\,\max\{\pi,1-\pi\}\le k \why{\(\DT\le\DT^{\rm N}\), equation~\eqref{eq:coexistence-quiet}}
\end{align*}
So class \(L\) strictly prefers abstention (Lemma~A.3(b)), and it cannot sell with positive probability. Hence
\(s_-=0\).
\end{derivation}

\paragraph{Step 2: only a purchase can recruit.} With \(s_-=0\), the sale has the posterior \(\mu_-(0)=\pi<\tau\).
For the zero order,
\begin{align*}
\frac{b_0+\lambda-s_+}{b_0+\lambda} &\le 1 \why{\(s_+\ge0\)}\\
\frac{\mu_0}{1-\mu_0}=\frac{\pi}{1-\pi}\cdot\frac{b_0+\lambda-s_+}{b_0+\lambda} &\le \frac{\pi}{1-\pi}
\why{multiply by the prior odds \(>0\)}
\end{align*}
The posterior rises with its odds, so \(\mu_0\le\pi<\tau\). So \(e(0)=e(-1)=\rho\). Only the posterior after a
purchase can reach the threshold.

\paragraph{Step 3: the purchase must recruit.} In an informative market class \(L\) abstains (Step~1), so class
\(H\) buys with positive probability: \(s_+>0\).

\begin{derivation}[two conditions for an informative market]
\emph{The purchase recruits.} Suppose \(\mu_+(s_+)<\tau\). Then \(e(1)=\rho\), and
\begin{align*}
K_H &= \rho\,\DT\,[1-\mu_+(s_+)] \why{\(e(1)=\rho\)}\\
&< \rho\,\DT\,(1-\pi) \why{\(\mu_+(s_+)>\pi\) for \(s_+>0\)}\\
&\le k \why{\(\DT\le\DT^{\rm N}\)}
\end{align*}
So class \(H\) would strictly prefer abstention, which contradicts \(s_+>0\). Hence \(\mu_+(s_+)\ge\tau\), and
\(e(1)=1\) under the default challenger rule.

\emph{The threshold is within reach.} The posterior \(\mu_+(s_+)\ge\tau\) at a rate \(s_+\le\lambda\) needs
\(\tau\le\mu_+(s_+)\le\mu_+(\lambda)=M\). So \(\tau\le M\).

\emph{The spread pays for a purchase at the threshold.} Class \(H\) buys with positive probability, so its purchase
is a best order: \(K_H\ge k\) (Lemma~A.3(b)). Then
\begin{align*}
k &\le K_H=e(1)\,\DT\,[1-\mu_+(s_+)] \why{Lemma~A.3(b)}\\
&\le \DT\,[1-\mu_+(s_+)] \why{\(e(1)\le1\)}\\
&\le \DT\,(1-\tau)=J \why{\(\mu_+(s_+)\ge\tau\)}
\end{align*}
Since \(\tau\le M<1\), \(1-\tau>0\), and \(k\le\DT(1-\tau)\) reads \(\DT\ge k/(1-\tau)\).
\end{derivation}

So, given \(\pi<\tau\) and \(\DT\le\DT^{\rm N}\), an informative market needs \(\tau\le M\) and
\(k/(1-\tau)\le\DT\).

\paragraph{Step 4: the purchase rate.} Let \(\tau\le M\) and \(J\ge k\). Write
\[
K(s)=\DT\,[1-\mu_+(s)]=\frac{\DT\,(1-\pi)\,b}{b+\pi s}
\]
for the return of a purchase at the rate \(s\) when it recruits. It falls strictly in \(s\). At \(s=s^*\) it is
\(\DT(1-\tau)=J\ge k\), and at \(s=\lambda\) it is \(\DT(1-M)\). Two cases follow.
\begin{itemize}
\item \emph{Case 1: \(\DT(1-M)\ge k\), that is \(\DT\ge\DT^{\rm F}\).} A full purchase still pays at the top
posterior. Class \(H\) buys for sure: \(s_+=\lambda\), the posterior is \(M\ge\tau\), and \(e(1)=1\) (at \(M=\tau\)
the indifferent challenger prepares). Its return is \(\DT(1-M)\ge k\). No other rate works. A rate below \(s^*\)
does not recruit (Step~3). A rate \(s\) in \([s^*,\lambda)\) has \(K(s)>K(\lambda)\ge k\), so class \(H\) strictly
prefers to buy and cannot mix.
\item \emph{Case 2: \(\DT(1-M)<k\le J\).} Then \(1-M<1-\tau\), so \(\tau<M\). The full purchase fails, since class
\(H\) would strictly prefer to abstain at \(K(\lambda)<k\). The function \(K\) is continuous and strictly decreasing
on \([s^*,\lambda]\), from \(J\ge k\) down to \(\DT(1-M)<k\). By the intermediate value theorem
(Section~\ref{sec:toolkit-monotone}), exactly one rate \(s_+\in[s^*,\lambda)\) has \(K(s_+)=k\). At this rate class
\(H\) is indifferent and mixes. If \(J>k\), then \(s_+>s^*\) and the posterior lies strictly above \(\tau\): the paper
calls this a \emph{mixed purchase that recruits} (label R). If \(J=k\), then \(s_+=s^*\) and the posterior equals
\(\tau\). The challenger is then indifferent and, under the default rule, prepares: the paper calls this the
\emph{threshold pair} (label T).
\end{itemize}

\begin{intuition}[why class \(H\) mixes]
The return of a purchase falls as class \(H\) buys more often, because a purchase then reveals more and leaves less
doubt to earn on. When even the full purchase leaves too little doubt to cover \(k\), class \(H\) buys only so often
that the doubt left after a purchase pays exactly \(k\).
\end{intuition}

\subsection{The purchase probability and entry}

While class \(H\) mixes, the purchase recruits and earns exactly the trading cost:
\begin{equation}\label{eq:coexistence-indiff}
\DT\,[1-\mu(1)]=k,\qquad\text{so}\qquad \mu(1)=1-\frac{k}{\DT}.
\end{equation}
The posterior after a purchase is then fixed by the spread and the trading cost alone. Bayes' rule turns it into the
purchase rate.

\begin{derivation}[the purchase rate while class \(H\) mixes]
\begin{align*}
\frac{(1-\pi)\,b}{b+\pi s_+} &= \frac{k}{\DT} \why{\(1-\mu_+(s_+)\) and \eqref{eq:coexistence-indiff}}\\
b+\pi s_+ &= \frac{(1-\pi)\,b\,\DT}{k} \why{multiply by \((b+\pi s_+)\DT/k>0\)}\\
\pi s_+ &= \frac{(1-\pi)\,b\,\DT}{k}-b=b\Bigl[\frac{(1-\pi)\DT}{k}-1\Bigr] \why{subtract \(b\); factor \(b\)}\\
s_+ &= \frac{b}{\pi}\Bigl[\frac{(1-\pi)\DT}{k}-1\Bigr] \why{divide by \(\pi>0\)}
\end{align*}
The probability with which class \(H\) buys is \(\sigma_H(1)=s_+/\lambda\):
\begin{align*}
\sigma_H(1) &= \frac{b}{\pi\lambda}\Bigl[\frac{(1-\pi)\DT}{k}-1\Bigr] \why{divide by \(\lambda\)}\\
&= \frac{1}{\pi u}\Bigl[\frac{(1-\pi)\DT}{k}-1\Bigr] \why{\(b/\lambda=1/u\)}
\end{align*}
\end{derivation}

The rate reaches \(\lambda\) exactly at the spread \(\DT^{\rm F}\):
\begin{align*}
\frac{b}{\pi}\Bigl[\frac{(1-\pi)\DT}{k}-1\Bigr]\ge\lambda
&\iff \frac{(1-\pi)\DT}{k}-1\ge\frac{\pi\lambda}{b} \why{multiply by \(\pi/b>0\)}\\
&\iff \frac{(1-\pi)\DT}{k}\ge\frac{b+\pi\lambda}{b} \why{add \(1=b/b\)}\\
&\iff \DT\,\frac{(1-\pi)\,b}{b+\pi\lambda}\ge k \why{multiply by \(kb/(b+\pi\lambda)>0\)}\\
&\iff \DT\,(1-M)\ge k \why{\(1-M=(1-\pi)b/(b+\pi\lambda)\)}
\end{align*}
So the two cases of Step~4 fit in one formula, the rate of Appendix~A.4:
\begin{equation}\label{eq:coexistence-rate}
s_+=\min\Bigl\{\lambda,\;\frac{b}{\pi}\Bigl[\frac{(1-\pi)\DT}{k}-1\Bigr]\Bigr\}.
\end{equation}

Entry is the probability that the challenger prepares. It adds up, over the three orders, the probability of the
order times entry after it.

\begin{derivation}[entry in the informative market]
The probability of a purchase is
\begin{align*}
\Pr(X=1) &= \pi f_H(1)+(1-\pi)f_L(1) \why{total probability over the classes}\\
&= \pi(b+s_+)+(1-\pi)\,b=b+\pi s_+ \why{\(\pi b+(1-\pi)b=b\)}
\end{align*}
With \(e(0)=e(-1)=\rho\) and \(e(1)=1\),
\begin{align*}
\Ent &= \sum_x\Pr(X=x)\,e(x) \why{total probability over the orders}\\
&= \rho\,[1-\Pr(X=1)]+1\cdot\Pr(X=1) \why{entry \(\rho\) after the two other orders}\\
&= \rho+(1-\rho)\Pr(X=1)=\rho+(1-\rho)(b+\pi s_+) \why{collect terms}
\end{align*}
While class \(H\) mixes, line 2 of the rate derivation gives \(b+\pi s_+=(1-\pi)b\DT/k\). When class \(H\) buys for
sure, \(b+\pi s_+=b+\pi\lambda\). The map \(s\mapsto b+\pi s\) rises, so \(b+\pi\min\{\lambda,\tilde s\}=\min\{b+\pi
\lambda,\,b+\pi\tilde s\}\), and
\begin{equation}\label{eq:coexistence-entry}
\Ent=\rho+(1-\rho)\min\Bigl\{\frac{(1-\pi)\,b\,\DT}{k},\;b+\pi\lambda\Bigr\}.
\end{equation}
\end{derivation}

This is the entry formula of Proposition~3(a), equation~(A.8) of the paper. When class \(H\) buys for sure it equals
\eqref{eq:ccc-entry}, entry against the strong incumbent of Proposition~2. The two terms of the minimum are equal
exactly when \((1-\pi)b\DT/k=b+\pi\lambda\), that is, when \(\DT(1-\pi)b/(b+\pi\lambda)=k\), that is, at
\(\DT=\DT^{\rm F}\). So:
\begin{itemize}
\item for \(\DT<\DT^{\rm F}\), entry is \(\rho+(1-\rho)(1-\pi)b\DT/k\), a straight line in \(\DT\) with the positive
slope \((1-\rho)(1-\pi)b/k\): it rises strictly with the spread;
\item for \(\DT\ge\DT^{\rm F}\), entry is the constant \(\rho+(1-\rho)(b+\pi\lambda)\);
\item at \(\DT^{\rm F}\) the two pieces meet, so entry has no jump.
\end{itemize}
The prior, \(b\), \(\lambda\), \(k\) and \(\rho\) are fixed, so entry depends on the incumbent only through the
spread. The threshold decides only whether the economy lies in the region. Between two incumbents at which the
markets coexist, entry in the informative market is at least as high against the one with the wider spread.

\begin{intuition}[why entry rises with the spread]
While class \(H\) mixes, a purchase that recruits earns exactly \(k\). A wider spread raises the return per unit of
doubt, so the purchase must leave less doubt. The purchase rate rises, purchases become more frequent, and each
purchase recruits. Entry rises with them until class \(H\) buys for sure.
\end{intuition}

There is a second way to see the entry formula. It is the guide's own.
\begin{align*}
\Pr(X=1)\,[1-\mu(1)] &= \Pr(X=1\text{ and }j=L) \why{\(1-\mu(1)=\Pr(j=L\mid X=1)\)}\\
&= (1-\pi)\,b \why{in class \(L\) only a liquidity trader buys}
\end{align*}
So \(\Pr(X=1)=(1-\pi)b/[1-\mu(1)]\). While class \(H\) mixes, \(1-\mu(1)=k/\DT\), so \(\Pr(X=1)=(1-\pi)b\DT/k\). In
words, the class-\(L\) purchases are a fixed pool. Indifference fixes the share of class \(L\) among all purchases at
\(k/\DT\). A wider spread needs a smaller share, so more class-\(H\) purchases must join the fixed pool.

Test the formula at its two ends.
\begin{itemize}
\item \emph{The left edge, \(J=k\).} Then \(\DT=k/(1-\tau)\), the posterior after a purchase is \(1-k/\DT=\tau\), and
the rate is \(s^*\): the threshold pair. Entry is \(\rho+(1-\rho)(1-\pi)b/(1-\tau)\), from
\((1-\pi)b\DT/k\) at \(\DT=k/(1-\tau)\).
\item \emph{The spread \(\DT^{\rm F}\).} Then \(\sigma_H(1)=1\), and entry reaches \(\rho+(1-\rho)(b+\pi\lambda)\).
\end{itemize}

\begin{watchout}
\begin{itemize}
\item A mixed purchase needs the default investor rule. If an indifferent investor abstains, class \(H\) cannot mix,
and at spreads below \(\DT^{\rm F}\) no informative market exists next to the quiet market.
\item The threshold pair at \(J=k\) needs the default challenger rule: the posterior equals \(\tau\), and an
indifferent high-cost challenger must prepare.
\item The entry formula describes one informative market at one incumbent. A statement across incumbents compares
two equilibria, each solved again in its own economy.
\end{itemize}
\end{watchout}

\begin{checkyourself}\label{q:coexistence-mix}
Let \(\pi=1/2\). Write the purchase probability \(\sigma_H(1)\) of class \(H\), while it mixes, in terms of \(\DT\),
\(k\) and \(u\). Show that it rises strictly with \(\DT\) and equals one exactly when \(\DT=k(2+u)\). Check that
\(k(2+u)=\DT^{\rm F}\) at this prior.
\end{checkyourself}

% =====================================================================
\section{Proposition 3}\label{sec:coexistence-statement}

Here is the result as Section~4.3 of the paper states it.

\begin{paperresult}{Proposition 3 (Quiet and informative markets)}
Let the classes satisfy the class ranking~(1), let \(0<\pi<1\), \(0<\rho<1\), \(k>0\) and a slot \((\lambda,\beta)\)
be given, and let an indifferent investor mix over its best orders and an indifferent high-cost challenger prepare.
Take an incumbent distribution with a positive spread \(\DT\) and \(g_H>g_L\), and write \(\tau\) for its preparation
threshold.
\begin{enumerate}[label=(\alph*),leftmargin=2em]
\item A quiet market and an informative market coexist if and only if
\[
\pi<\tau\le M(\pi),\qquad \frac{k}{1-\tau}\le\DT\le\frac{k}{\rho\max\{\pi,1-\pi\}}.
\]
The informative market is then unique: class \(L\) abstains, only a purchase recruits a high-cost challenger, and
entry is
\[
\rho+(1-\rho)\min\Bigl\{\frac{(1-\pi)\,b\,\DT}{k},\;b+\pi\lambda\Bigr\},
\]
which rises strictly with the spread while class \(H\) mixes between a purchase and abstention and is constant once
class \(H\) buys for sure.
\item Along a path of such distributions, the strengths at which a quiet market coexists with an informative market
in which class \(H\) buys for sure form an interval, and every other strength with coexistence lies below it. Every
gap in coexistence lies at spreads too narrow for a sure purchase to pay.
\item Along a path on which the product of the spread and \(g_H-c_H\) does not fall with strength, the strengths
with coexistence form an interval, at every prior. Uniform incumbents on \([0,r]\), with \(r\) in a closed interval
above the reserve, satisfy this condition when every low-class value lies between the reserve and the lowest
strength, every high-class value is at least the highest strength, and class \(H\)'s mean value exceeds the highest
strength by at least \(c_H\).
\end{enumerate}
\end{paperresult}

A plain reading, line by line:
\begin{itemize}
\item \emph{Hypotheses.} The classes are ranked. The prior, the floor \(\rho\), the trading cost and the slot are
free; Appendix~A.4 fixes \(0<\lambda<1\) and \(0<\beta<1/2\). The tie rules are the default ones. The incumbent law
has a positive spread and \(g_H>g_L\), so the price reveals the posterior and the threshold \(\tau\) exists; the
appendix also takes the law continuous. Nothing like (H1) to (H3) is assumed.
\item \emph{Part (a).} The four conditions of Section~\ref{sec:coexistence-words} decide coexistence, in both
directions. In the region the informative market is unique, and it is the market of
Section~\ref{sec:coexistence-informative}.
\item \emph{Part (b).} Move along a path of incumbents that grow stronger. The strengths at which the informative
market has a full purchase form one interval, and the strengths with a mixed purchase lie to its left. A hole in the
coexistence set can occur only where the spread is below \(\DT^{\rm F}\).
\item \emph{Part (c).} A condition on the path, the \emph{balance condition}, keeps the whole coexistence set in one
piece. The last sentence gives a family of uniform incumbents that meets it.
\end{itemize}

Appendix~A.4 restates the proposition at more length. Its part (a) adds the purchase rate
\eqref{eq:coexistence-rate}; the type of the informative market (a full purchase if \(\DT\ge\DT^{\rm F}\), the
threshold pair if \(\DT<\DT^{\rm F}\) and \(J=k\), a mixed purchase that recruits if \(\DT<\DT^{\rm F}\) and
\(J>k\)); the fact that the region holds exactly two equilibria; the fact that entry depends on the incumbent only
through the spread; and the shape of the region. Its part (c) has four items, and a part (d) shows that an
equilibrium exists at the prior \(1/2\). The guide covers each item where it fits.

\begin{watchout}
Proposition~3 uses the default tie rules. Proposition~2 holds under all four tie rules, because its strict
conditions keep every player away from a tie. Here ties are part of the result: a mixed purchase is an indifference
of the investor, the threshold pair is an indifference of the challenger, and the edge \(\DT=\DT^{\rm N}\) is an
indifference of the investor in the quiet market. Section~\ref{sec:coexistence-ties} says what changes under the other
rules.
\end{watchout}

% =====================================================================
\section{The proof of part (a)}\label{sec:coexistence-proof-a}

The proof in Appendix~A.4 has four steps. Step~1 rules out a prior at or above the threshold. Three lemmas then
describe every equilibrium when \(\pi<\tau\), one piece at a time. Steps~2 to~4 apply them. Fix an incumbent law
\(F\) and the default tie rules.

\paragraph{Step 1: a prior at or above the threshold.} Let \(\pi\ge\tau\).

\begin{derivation}[no coexistence when \(\pi\ge\tau\)]
\emph{The quiet market.} Every posterior is \(\pi\ge\tau\), so entry is one after every order (Lemma~A.2(d); at
\(\pi=\tau\) the indifferent challenger prepares). The deviations earn \(K_H=\DT(1-\pi)\) and \(K_L=\DT\pi\). So a
quiet market needs \(\DT\max\{\pi,1-\pi\}\le k\).

\emph{No informative market under this inequality.} Take a candidate with informed trade. By Lemma~A.3(a) it has no
wrong-signed order. If \(s_+>0\), then \(\mu(1)=\mu_+(s_+)>\pi\ge\tau\), so \(e(1)=1\), and
\begin{align*}
K_H &= \DT\,[1-\mu_+(s_+)] \why{\(e(1)=1\)}\\
&< \DT\,(1-\pi)\le k \why{\(\mu_+(s_+)>\pi\); the quiet-market inequality}
\end{align*}
So class \(H\) strictly prefers abstention. If \(s_->0\), then \(\mu(-1)=\mu_-(s_-)<\pi\) and \(e(-1)\le1\), so
\begin{align*}
K_L &= e(-1)\,\DT\,\mu_-(s_-) \why{equation~\eqref{eq:trading-returns}}\\
&\le \DT\,\mu_-(s_-)<\DT\,\pi\le k \why{\(e(-1)\le1\); \(\mu_-(s_-)<\pi\); the inequality}
\end{align*}
So class \(L\) strictly prefers abstention. In both cases the candidate fails (E5).
\end{derivation}

So coexistence needs \(\pi<\tau\). In words: when the prior already recruits, information cannot add entry. It can
only take doubt away, so every informed order earns less than a deviation from the quiet market.

\subsection{Three lemmas}

From here on let \(\pi<\tau\). Appendix~A.4 states three lemmas. They hold under each pair of tie rules, and they do
not assume that the quiet market is an equilibrium.

\begin{paperresult}{Lemma A.7 (The high class)}
Let \(\pi<\tau\). In a candidate with no wrong-signed order, the pair \((s_+,e(1))\) satisfies (E1) to (E3) at the
order \(1\), and (E5) for class \(H\), exactly when it is one of these five pairs.
\begin{itemize}[leftmargin=2.5em]
\item[(N)] \(s_+=0\) and \(e(1)=\rho\), with \(\rho\DT(1-\pi)\le k\).
\item[(Q)] \(0<s_+<\lambda\), \(\mu_+(s_+)<\tau\), \(e(1)=\rho\) and \(\rho\DT[1-\mu_+(s_+)]=k\). This pair needs an
indifferent investor who may mix.
\item[(R)] \(0<s_+<\lambda\), \(\mu_+(s_+)>\tau\), \(e(1)=1\) and \(\DT[1-\mu_+(s_+)]=k\). This pair needs an
indifferent investor who may mix.
\item[(T)] \(0<s_+<\lambda\), \(\mu_+(s_+)=\tau\), that is \(s_+=s^*\), and \(e(1)=e^*:=k/J\). When an indifferent
high-cost challenger prepares, this requires \(e^*=1\), that is \(J=k\). When an indifferent high-cost challenger may
prepare with any probability, it requires \(\rho\le e^*\le1\), that is \(\rho J\le k\le J\). This pair needs an
indifferent investor who may mix.
\item[(F)] \(s_+=\lambda\). Here \(e(1)=1\) if \(M>\tau\) and \(e(1)=\rho\) if \(M<\tau\). If \(M=\tau\), then
\(e(1)=1\) when an indifferent high-cost challenger prepares, and \(e(1)\in[\rho,1]\) when it may prepare with any
probability. The return satisfies \(e(1)\DT(1-M)\ge k\) when an indifferent investor may mix, and
\(e(1)\DT(1-M)>k\) when an indifferent investor abstains.
\end{itemize}
Each equation in (Q) and (R) has at most one solution. If \(\tau\ge M\), and in particular if \(\tau\ge1\), the pairs
(R) and (T) do not occur, and (F) has \(e(1)=\rho\) when \(\tau>M\).
\end{paperresult}

A plain reading. The split by class (Section~\ref{sec:trading-split}) lets the lemma look at class \(H\) and the
purchase alone. It lists every purchase rate, with the entry after a purchase, that agrees with Bayes' rule, the
challenger's choice and class \(H\)'s best order. The five letters name the five shapes: no purchase (N); a mixed
purchase that does not recruit, at the entry floor (Q); a mixed purchase that recruits (R); a mixed purchase that
lands exactly on the threshold (T); and a full purchase (F).

\begin{derivation}[proof of Lemma A.7]
By Lemma~A.3(b), (E5) for class \(H\) reads as follows.
\begin{itemize}
\item At \(s_+=0\), class \(H\) never buys, so abstention must be best: \(K_H\le k\).
\item At \(s_+=\lambda\), class \(H\) always buys, so the purchase must be best: \(K_H\ge k\) when an indifferent
investor may mix, and \(K_H>k\) when it abstains, because that rule picks abstention at a tie.
\item At \(0<s_+<\lambda\), class \(H\) mixes, so it must be indifferent: \(K_H=k\). A mix needs an indifferent
investor who may mix.
\end{itemize}
Here \(K_H=e(1)\DT[1-\mu_+(s_+)]\), and \(e(1)\) follows from Lemma~A.2(d) at the posterior \(\mu_+(s_+)\). Go through
the cases.
\begin{itemize}
\item \(s_+=0\): \(\mu_+(0)=\pi<\tau\), so \(e(1)=\rho\) and \(K_H=\rho\DT(1-\pi)\le k\). This is (N).
\item \(0<s_+<\lambda\) with \(\mu_+(s_+)<\tau\): \(e(1)=\rho\), and \(K_H=k\) is the equation of (Q).
\item \(0<s_+<\lambda\) with \(\mu_+(s_+)>\tau\): \(e(1)=1\), and \(K_H=k\) is the equation of (R).
\item \(0<s_+<\lambda\) with \(\mu_+(s_+)=\tau\): then \(s_+=s^*<\lambda\), so \(\pi<\tau<M<1\) and \(J=\DT(1-\tau)>0\).
The condition \(K_H=e(1)\DT(1-\tau)=e(1)J=k\) gives \(e(1)=k/J=e^*\). If the indifferent challenger prepares,
\(e(1)=1\), so \(J=k\). If it may prepare with any probability, \(e(1)\) can be any value in \([\rho,1]\), so the
condition is \(\rho\le k/J\le1\). This is (T).
\item \(s_+=\lambda\): \(\mu_+(\lambda)=M\). Lemma~A.2(d) at \(M\) gives \(e(1)\), and the condition on \(K_H\) is that
of (F).
\end{itemize}
The function \(1-\mu_+\) falls strictly and \(\DT>0\), so \(\rho\DT[1-\mu_+(s)]\) and \(\DT[1-\mu_+(s)]\) fall
strictly in \(s\). Each equals \(k\) at most once. If \(\tau\ge M\), then \(\mu_+(s_+)<M\le\tau\) for every
\(s_+<\lambda\), so (R) and (T) do not occur; and \(M<\tau\) gives \(e(1)=\rho\) in (F).
\end{derivation}

\begin{paperresult}{Lemma A.8 (The low class)}
Let \(\pi<\tau\). Then \(\mu_-(s_-)\le\pi<\tau\) for every \(s_-\), so \(e(-1)=\rho\) in every candidate. In a
candidate with no wrong-signed order, the rate \(s_-\) satisfies (E1) to (E3) at the order \(-1\), and (E5) for class
\(L\), exactly when it is one of these:
\begin{itemize}[leftmargin=2.5em]
\item[(N)] \(s_-=0\), with \(\rho\DT\pi\le k\);
\item[(Q)] \(0<s_-<\lambda\), with \(\rho\DT\mu_-(s_-)=k\); this rate needs an indifferent investor who may mix;
\item[(F)] \(s_-=\lambda\), with \(\rho\DT m\ge k\) when an indifferent investor may mix, and \(\rho\DT m>k\) when an
indifferent investor abstains.
\end{itemize}
When an indifferent investor may mix, exactly one rate qualifies:
\[
\hat s_-=\begin{cases}0, & \rho\DT\pi\le k,\\[2pt]
\dfrac{b\,(\pi\rho\DT/k-1)}{1-\pi}, & \rho\DT m<k<\rho\DT\pi,\\[6pt]
\lambda, & \rho\DT m\ge k.\end{cases}
\]
When an indifferent investor abstains, no rate qualifies if \(\rho\DT m\le k<\rho\DT\pi\).
\end{paperresult}

A plain reading. When the prior lies below the threshold, a sale never recruits, so class \(L\) always trades at
floor entry. Its sale rate is set by one comparison with \(k\). The more often class \(L\) sells, the more a sale
reveals and the less it earns. So class \(L\) abstains when even a sale at the prior does not pay, sells for sure when
even the full sale pays, and mixes in between.

\begin{derivation}[proof of Lemma A.8]
Since \(\mu_-\) falls, \(\mu_-(s_-)\le\mu_-(0)=\pi<\tau\), so Lemma~A.2(d) gives \(e(-1)=\rho\). Then
\(K_L=\rho\DT\mu_-(s_-)\) is continuous and strictly decreasing in \(s_-\), from \(\rho\DT\pi\) at \(s_-=0\) to
\(\rho\DT m\) at \(s_-=\lambda\). By Lemma~A.3(b), (E5) for class \(L\) requires \(K_L\le k\) at \(s_-=0\);
\(K_L\ge k\) at \(s_-=\lambda\) when an indifferent investor may mix, and \(K_L>k\) when it abstains; and \(K_L=k\)
at \(0<s_-<\lambda\), which needs an indifferent investor who may mix. This gives (N), (Q) and (F).

\emph{One rate when the investor may mix.}
\begin{itemize}
\item If \(\rho\DT\pi\le k\): then \(K_L<\rho\DT\pi\le k\) for every \(s_->0\), so only \(s_-=0\) qualifies.
\item If \(\rho\DT m\ge k\): then \(K_L>\rho\DT m\ge k\) for every \(s_-<\lambda\), so only \(s_-=\lambda\) qualifies.
\item Otherwise \(\rho\DT m<k<\rho\DT\pi\). The intermediate value theorem and strict monotonicity give exactly one
interior root, and neither end qualifies.
\end{itemize}
The interior root solves \(\rho\DT\mu_-(s)=k\):
\begin{align*}
\frac{\pi b}{b+(1-\pi)s} &= \frac{k}{\rho\DT} \why{\(\mu_-(s)\); divide by \(\rho\DT\)}\\
b+(1-\pi)s &= \frac{\pi b\,\rho\DT}{k} \why{multiply by \([b+(1-\pi)s]\rho\DT/k\)}\\
(1-\pi)s &= b\Bigl(\frac{\pi\rho\DT}{k}-1\Bigr) \why{subtract \(b\); factor \(b\)}\\
s &= \frac{b\,(\pi\rho\DT/k-1)}{1-\pi} \why{divide by \(1-\pi>0\)}
\end{align*}

\emph{No rate when the investor abstains.} Let \(\rho\DT m\le k<\rho\DT\pi\). The rate \(s_-=0\) fails, because
\(K_L=\rho\DT\pi>k\) there, so class \(L\) strictly prefers to sell. The rate \(s_-=\lambda\) fails, because
\(K_L=\rho\DT m\le k\) there, and this rule needs a strict inequality. An interior rate needs a mix.
\end{derivation}

\begin{paperresult}{Lemma A.9 (Product structure)}
Let \(\pi<\tau\), and take rates with \(s_+\ge s_-\). Then \(\mu_0(s_+,s_-)\le\pi<\tau\). Under each pair of tie rules,
the tuple with the rates \((s_+,s_-)\) is an equilibrium if and only if these four conditions hold:
\begin{enumerate}[label=(\roman*),leftmargin=2.5em]
\item \((s_+,e(1))\) is a pair of Lemma~A.7;
\item \(s_-\) is a rate of Lemma~A.8;
\item \(e(0)=e(-1)=\rho\);
\item the prices are \(P(x)=t_0+e(x)\,w(\mu(x))\), and \(\bar\mu(P(x))=\mu(x)\).
\end{enumerate}
\end{paperresult}

A plain reading. An equilibrium is built from one piece for class \(H\) and one piece for class \(L\), with entry at
the floor after the zero order and a sale, and prices from the price formula. The rule ``purchase rate at least sale
rate'' keeps the zero order below the threshold, so the two pieces do not interact. Recall
\(w(\mu)=t_L-t_0+\DT\mu\) from Section~\ref{sec:beliefs-priceformula}.

\begin{derivation}[proof of Lemma A.9]
\emph{The zero order.} If \(s_+\ge s_-\), then \(b_0+\lambda-s_+\le b_0+\lambda-s_-\). So the odds of \(\mu_0\) are
at most \(\pi/(1-\pi)\), and \(\mu_0\le\pi<\tau\).

\emph{Necessity.} Take an equilibrium. The posteriors \(\mu_0\) and \(\mu_-(s_-)\) lie below \(\tau\), so
Lemma~A.2(d) gives \(e(0)=e(-1)=\rho\): condition (iii). Lemma~A.2(b) and~(c) give the prices and the price beliefs:
condition (iv). By Lemma~A.3(d), (E5) for class \(H\) involves only \((s_+,\mu(1),e(1))\), and (E5) for class \(L\)
only \((s_-,\mu(-1),e(-1))\). So (E5) gives (i) and (ii).

\emph{Sufficiency.} Take rates and entry that satisfy (i) to (iv). Check each condition of the definition.
\begin{itemize}
\item (E1): the posteriors are computed from \(f_H\) and \(f_L\) by Bayes' rule.
\item (E2): entry lies in \([\rho,1]\) and does not fall when the posterior rises. It is \(\rho\) at \(\mu_-(s_-)\) and
\(\mu_0\), both at most \(\pi<\tau\), and it is \(e(1)\ge\rho\) at \(\mu_+(s_+)\ge\pi\), with \(e(1)=1\) when
\(\mu_+(s_+)>\tau\). Next, two orders with a common posterior. One of them is \(0\) or \(-1\), so the common
posterior is at most \(\pi\). If the order \(1\) is one of them, then \(\mu_+(s_+)\le\pi\) forces \(s_+=0\), the
pair (N), with \(e(1)=\rho\). So orders with a common posterior have a common entry. By the result on prices built
from a profile (Section~\ref{sec:beliefs-constructed}), orders share a price exactly when they share the posterior
and entry, the rule \(\zeta(P(x))=[e(x)-\rho]/(1-\rho)\) is a function of the price, and the belief at each price
is the common posterior of its orders. This is (E2).
\item (E3): at each price the belief is the posterior of its orders, and the stated entry is the one that
Lemma~A.2(d) allows at that posterior, under the challenger rule in use. Condition (i) states it after a purchase;
the other two posteriors lie below \(\tau\).
\item (E4): the price formula \eqref{eq:beliefs-price} with the stated entry.
\item (E5): by Lemma~A.3(d), class \(H\)'s payoffs involve only \((s_+,\mu(1),e(1))\), and class \(L\)'s only
\((s_-,\mu(-1),e(-1))\). Conditions (i) and (ii) state optimality for each, and Lemma~A.3(a) rules out the
wrong-signed orders.
\item (E6): truthful bids are weakly dominant (Chapter~\ref{ch:auction}).
\end{itemize}
\end{derivation}

\begin{watchout}
The condition \(s_+\ge s_-\) does real work. It puts the posterior after the zero order at or below \(\pi\), below the
threshold. Without it, class \(L\) could sell more often than class \(H\) buys, the zero order would become good news
about class \(H\), and its posterior could reach \(\tau\). This direct bound also explains why the chapter never needs
\(\beta\le1/3\): the posterior bound at the zero order (Lemma~A.1(c)) is not used.
\end{watchout}

\subsection{Steps 2 to 4}

\paragraph{Step 2: the quiet market when \(\pi<\tau\).} A quiet market has \(s_+=s_-=0\), so \(s_+\ge s_-\) and
Lemma~A.9 applies. It is an equilibrium exactly when \((0,e(1))\) is the pair (N) of Lemma~A.7 and \(0\) is the rate
(N) of Lemma~A.8:
\[
\rho\DT(1-\pi)\le k\quad\text{and}\quad\rho\DT\pi\le k,\qquad\text{that is}\qquad\DT\le\DT^{\rm N}.
\]
This is \eqref{eq:coexistence-quiet}. Its entry is \(\rho\) after every order, and condition (iv) fixes its prices,
so there is exactly one quiet market.

\paragraph{Step 3: informative markets when \(\pi<\tau\) and \(\DT\le\DT^{\rm N}\).}

\begin{derivation}[necessity]
Since \(\rho\DT\pi\le k\), Lemma~A.8 gives \(s_-=0\) in every equilibrium. Then \(s_+\ge0=s_-\), and by Lemma~A.9 the
equilibria are the pairs of Lemma~A.7 with \(s_-=0\) and entry \(\rho\) after the orders \(0\) and \(-1\). An
informative market has \(s_+>0\).
\begin{itemize}
\item The pair (Q) is impossible: \(\rho\DT[1-\mu_+(s_+)]<\rho\DT(1-\pi)\le k\) for \(s_+>0\).
\item The pair (F) with \(M<\tau\) is impossible for the same reason: its return is \(\rho\DT(1-M)<k\).
\item The other pairs, (R), (T) and (F) with \(M\ge\tau\), need \(\mu_+(s_+)\ge\tau\) for some \(s_+\le\lambda\), so
\(\tau\le M\). Since \(e(1)\le1\), they also need \(k\le e(1)\DT[1-\mu_+(s_+)]\le\DT(1-\tau)=J\), that is
\(\DT\ge k/(1-\tau)\).
\end{itemize}
With Step~1, coexistence implies \(\pi<\tau\le M\) and \(k/(1-\tau)\le\DT\le\DT^{\rm N}\).
\end{derivation}

\begin{derivation}[sufficiency and uniqueness]
Let the four conditions hold. Then \(\tau\le M<1\) and \(J\ge k\), and Step~2 gives the quiet market.
\begin{itemize}
\item If \(\DT\ge\DT^{\rm F}\), that is \(\DT(1-M)\ge k\), the pair (F) holds with \(e(1)=1\); at \(M=\tau\) the
indifferent challenger chooses this. It is the only informative pair: at every \(s_+<\lambda\) with
\(\mu_+(s_+)\ge\tau\), the return \(\DT[1-\mu_+(s_+)]\) exceeds \(\DT(1-M)\ge k\), so neither (R) nor (T) holds.
\item If \(\DT<\DT^{\rm F}\), then \(\DT(1-M)<k\le J=\DT(1-\tau)\), so \(\tau<M\) and (F) fails. The return
\(s\mapsto\DT[1-\mu_+(s)]\) is continuous and strictly decreasing on \([s^*,\lambda]\), from \(J\ge k\) to
\(\DT(1-M)<k\). So it equals \(k\) at exactly one \(s_+\in[s^*,\lambda)\). At \(s_+=s^*\), which happens exactly when
\(J=k\), this is (T) with \(e^*=1\). At \(s_+>s^*\) it is (R). No other informative pair exists.
\end{itemize}
In every case the low class abstains, entry after the orders \(0\) and \(-1\) is \(\rho\), and Lemma~A.9 completes the
tuple. The rate is \eqref{eq:coexistence-rate}, derived in Section~\ref{sec:coexistence-informative}.
\end{derivation}

So in the region the equilibria are exactly two: the quiet market and one informative market.

\paragraph{Step 4: entry and the shape of the region.} Section~\ref{sec:coexistence-informative} derived the entry
formula \eqref{eq:coexistence-entry} and its shape. Two facts about the region remain.

\begin{derivation}[the region is convex]
Since \(\tau\le M<1\), the inequality \(k/(1-\tau)\le\DT\) reads \(\tau\le1-k/\DT\). So the region is the set of
pairs \((\DT,\tau)\) with \(\DT>0\) and
\[
\pi<\tau,\qquad\tau\le M,\qquad\tau\le\varphi(\DT),\qquad\DT\le\DT^{\rm N},\qquad\text{where }\varphi(\DT)=1-\frac{k}{\DT}.
\]
The function \(\varphi\) is concave: \(\varphi'(\DT)=k/\DT^2\) falls as \(\DT\) rises. Take two pairs with
\(\tau_1\le\varphi(\Delta_1)\) and \(\tau_2\le\varphi(\Delta_2)\), and a weight \(a\in[0,1]\). Then
\begin{align*}
a\tau_1+(1-a)\tau_2 &\le a\,\varphi(\Delta_1)+(1-a)\,\varphi(\Delta_2) \why{each \(\tau_i\le\varphi(\Delta_i)\)}\\
&\le \varphi\bigl(a\Delta_1+(1-a)\Delta_2\bigr) \why{\(\varphi\) is concave}
\end{align*}
So the set below the graph of \(\varphi\) is convex: the segment between two of its points stays in it. Each of the
other three conditions defines a half-plane, which is convex. An intersection of convex sets is convex.
\end{derivation}

\begin{derivation}[the threshold cannot be too close to one]
Put the two spread bounds together:
\begin{align*}
\frac{k}{1-\tau} &\le \frac{k}{\rho\max\{\pi,1-\pi\}} \why{\(k/(1-\tau)\le\DT\le\DT^{\rm N}\)}\\
\frac{1}{1-\tau} &\le \frac{1}{\rho\max\{\pi,1-\pi\}} \why{divide by \(k>0\)}\\
\rho\max\{\pi,1-\pi\} &\le 1-\tau \why{reciprocals of positive numbers; the inequality turns}\\
\tau &\le 1-\rho\max\{\pi,1-\pi\} \why{rearrange}
\end{align*}
With \(\pi<\tau\) and \(\max\{\pi,1-\pi\}\ge\pi\):
\begin{align*}
\pi &< 1-\rho\max\{\pi,1-\pi\}\le1-\rho\pi \why{the line above}\\
\pi\,(1+\rho) &< 1 \why{add \(\rho\pi\)}\\
\pi &< \frac{1}{1+\rho} \why{divide by \(1+\rho>0\)}
\end{align*}
\end{derivation}

In words: a purchase that just reaches the threshold leaves the doubt \(1-\tau\) and must pay; an order against a
silent price leaves the doubt \(\max\{\pi,1-\pi\}\), at entry \(\rho\), and must not pay. So the doubt \(1-\tau\) must
be at least \(\rho\max\{\pi,1-\pi\}\). A prior that strongly favours class \(H\) leaves too little room between
itself and the largest admissible threshold, and coexistence needs \(\pi<1/(1+\rho)\).
Figure~\ref{fig:coexistence-region} draws the region.

\begin{figure}[htbp]
\centering
\begin{tikzpicture}[x=1cm,y=1cm,font=\small]
\fill[accent!15] plot[domain=0.2:0.37037,samples=40,smooth] ({10*\x},{6-1/\x}) -- (8,3.3) -- (8,1) -- cycle;
\draw[-{Stealth[length=2mm]}] (0,0) -- (10.6,0) node[right] {\(\DT\)};
\draw[-{Stealth[length=2mm]}] (0,0) -- (0,4.6) node[above] {\(\tau\)};
\draw[rulegrey,dotted] plot[domain=0.17:0.2,samples=10,smooth] ({10*\x},{6-1/\x});
\draw[rulegrey,dotted] plot[domain=0.37037:0.8,samples=40,smooth] ({10*\x},{6-1/\x});
\draw[accent,thick] plot[domain=0.2:0.37037,samples=40,smooth] ({10*\x},{6-1/\x});
\draw[accent,thick] (3.7037,3.3) -- (8,3.3);
\draw[accent,thick] (8,1) -- (8,3.3);
\draw[accent,thick,dotted] (2,1) -- (8,1);
\draw[dashed,ink] (3.7037,1) -- (3.7037,3.3);
\draw[dotted,rulegrey] (0,1) -- (2,1);
\draw[dotted,rulegrey] (0,3.3) -- (3.7037,3.3);
\draw[dotted,rulegrey] (2,0) -- (2,1);
\draw[dotted,rulegrey] (3.7037,0) -- (3.7037,1);
\draw[dotted,rulegrey] (8,0) -- (8,1);
\node[left] at (0,1) {\(\pi\)};
\node[left] at (0,3.3) {\(M(\pi)\)};
\node[below] at (2,0) {\(\frac{k}{1-\pi}\)};
\node[below] at (3.7037,0) {\(\DT^{\rm F}\)};
\node[below] at (8,0) {\(\DT^{\rm N}\)};
\node at (3.2,1.7) {R};
\node at (5.9,2.1) {F};
\node[left] at (2.75,2.6) {\(J=k\)};
\node[above] at (5.85,3.3) {\(\tau=M(\pi)\)};
\node[right] at (8,2.15) {\(\DT=\DT^{\rm N}\)};
\node[text=notegrey] at (5.6,0.6) {\(\tau=\pi\), not in the region};
\end{tikzpicture}
\caption{A picture to keep in mind: the coexistence region of Proposition~3(a) in the plane of the spread \(\DT\)
and the threshold \(\tau\). The shaded set holds the pairs with \(\pi<\tau\le M(\pi)\) and
\(k/(1-\tau)\le\DT\le\DT^{\rm N}\). Its left edge is the curve \(\tau=1-k/\DT\), where \(J=k\); this curve meets the
top edge at \(\DT^{\rm F}\). The dashed line divides the mixed purchase that recruits (R) from the full purchase (F).
Solid edges belong to the region, and the dotted lower edge does not. The positions are only a picture.}
\label{fig:coexistence-region}
\end{figure}

The left edge meets the top edge exactly at \(\DT^{\rm F}\): at \(\tau=M\), the condition \(\DT=k/(1-\tau)\) reads
\(\DT=k/(1-M)=\DT^{\rm F}\). So the part with a mixed purchase lies to the left of the dashed line, below the top edge
and right of the curve.

\begin{checkyourself}\label{q:coexistence-edge}
Take an incumbent on the left edge of the region, \(J=k\), with \(\pi<\tau<M\) and \(\DT\le\DT^{\rm N}\). Show that
the purchase rate \eqref{eq:coexistence-rate} equals \(s^*\). Explain why this informative market needs the default
challenger rule, and what entry after a purchase would have to be under the other challenger rule.
\end{checkyourself}

% =====================================================================
\section{Part (b): the strengths with a sure purchase}\label{sec:coexistence-proof-b}

Recall that a path is a family \((F_r)\) of such laws over an interval of strengths, with \(F_{r'}\) a first-order
stochastic strengthening of \(F_r\) for \(r'>r\) and every other primitive fixed (Section~\ref{sec:trading-path}).
Along a path, \(\DT(r)\) never falls, and \(g_H(r)\), \(g_L(r)\) and \(\gamma(r)=g_H(r)-g_L(r)\) never rise
(Lemma~A.5). Each strength is a separate economy, and the coexistence set at that strength comes from solving that
economy.

First the idea in words. As the incumbent strengthens, three of the four conditions of the region move in one
direction only. The condition \(\pi<\tau\), once true, stays true; the conditions \(\tau\le M\) and
\(\DT\le\DT^{\rm N}\), once false, stay false. The fourth condition, \(J\ge k\), has no fixed direction. For a sure
purchase, the fourth condition is replaced by \(\DT\ge\DT^{\rm F}\), which also moves in one direction. So the sure
purchase lives on one interval.

Two words from Appendix~A.4. A set of strengths is an \emph{up-set} if it contains every strength above each of its
members, and a \emph{down-set} if it contains every strength below each of its members.

\begin{derivation}[up-sets and down-sets meet in an interval]
Let \(S\) be an intersection of up-sets and down-sets, and let \(r'<r<r''\) with \(r'\) and \(r''\) in \(S\).
\begin{itemize}
\item Each up-set contains \(r'\), so it contains every strength above \(r'\), and in particular \(r\).
\item Each down-set contains \(r''\), so it contains every strength below \(r''\), and in particular \(r\).
\end{itemize}
So \(r\in S\). A set that contains every strength between two of its members is an interval.
\end{derivation}

\begin{derivation}[the direction of each condition]
\emph{Gross profit at a fixed belief never rises.} For every \(\mu\in[0,1]\),
\begin{align*}
B_r(\mu) &= \mu\,g_H(r)+(1-\mu)\,g_L(r) \why{definition of \(B\)}\\
&\ge \mu\,g_H(r')+(1-\mu)\,g_L(r')=B_{r'}(\mu)\quad(r'>r) \why{each \(g_j\) never rises; weights \(\ge0\)}
\end{align*}

\emph{The threshold conditions.} Since \(\gamma(r)>0\), \(B_r(\mu)-c_H=\gamma(r)\,[\mu-\tau(r)]\)
(Section~\ref{sec:trading-twonumbers}), so \(B_r(\mu)-c_H\) has the sign of \(\mu-\tau(r)\). Hence
\begin{align*}
\pi<\tau(r) &\iff B_r(\pi)<c_H, \why{sign rule at \(\mu=\pi\)}\\
\tau(r)\le M &\iff B_r(M)\ge c_H. \why{sign rule at \(\mu=M\)}
\end{align*}
If \(B_r(\pi)<c_H\) and \(r'>r\), then \(B_{r'}(\pi)\le B_r(\pi)<c_H\): the first set is an up-set. If
\(B_{r''}(M)\ge c_H\) and \(r<r''\), then \(B_r(M)\ge B_{r''}(M)\ge c_H\): the second set is a down-set.

\emph{The spread conditions.} \(\DT\) never falls, so \(\{r:\DT(r)\ge\DT^{\rm F}\}\) is an up-set and
\(\{r:\DT(r)\le\DT^{\rm N}\}\) is a down-set.
\end{derivation}

This argument uses gross profit at a fixed belief, not the threshold itself. So it needs no condition like
\(g_L\le c_H\), which the statement ``a stronger incumbent never lowers the threshold'' needs
(Section~\ref{sec:trading-path}).

\begin{derivation}[proof of part (b)]
\emph{The sure-purchase set is an interval.} By part (a), a quiet market coexists with a full purchase exactly where
the region holds and \(\DT(r)\ge\DT^{\rm F}\). There the condition \(J\ge k\) holds by itself:
\begin{align*}
J &= \DT(1-\tau) \why{definition}\\
&\ge \DT(1-M) \why{\(\tau\le M\)}\\
&\ge k \why{\(\DT\ge\DT^{\rm F}\)}
\end{align*}
So this set is \(S_{\rm F}=\{r:\pi<\tau(r)\le M,\ \DT^{\rm F}\le\DT(r)\le\DT^{\rm N}\}\): two up-sets and two
down-sets. It is an interval.

\emph{Every other strength with coexistence lies below it.} At a strength \(r\) with coexistence and a mixed purchase
that recruits, or the threshold pair, part (a) gives \(\DT(r)<\DT^{\rm F}\). Take any \(r'\in S_{\rm F}\). If \(r\ge r'\),
then \(\DT(r)\ge\DT(r')\ge\DT^{\rm F}\), because \(\DT\) never falls. That is false, so \(r<r'\).

\emph{A gap lies at a narrow spread.} The set \(A=\{r:\pi<\tau(r)\le M,\ \DT(r)\le\DT^{\rm N}\}\) is an interval (one
up-set and two down-sets), and it contains every strength with coexistence. Let coexistence hold at \(r'\) and \(r''\)
but not at \(r\), with \(r'<r<r''\). Then \(r\in A\). If \(\DT(r)\ge\DT^{\rm F}\), then \(r\in S_{\rm F}\), and
coexistence would hold at \(r\). So \(\DT(r)<\DT^{\rm F}\).
\end{derivation}

\begin{intuition}[why the sure purchase is never split]
Each condition that defines the sure-purchase stretch holds on one side of a single strength. A mixed purchase needs
a narrower spread than a sure one, and the spread only widens along the path. So the strengths with a mixed purchase
come first, and the sure-purchase stretch comes after them, in one piece. A hole can appear only where the economy
relies on a mixed purchase, that is, where \(J\ge k\) must hold by itself.
\end{intuition}

% =====================================================================
\section{Part (c): the balance condition}\label{sec:coexistence-proof-c}

The only condition of the region with no fixed direction is \(J\ge k\). Part (c) gives it a direction. Recall from
Section~\ref{sec:trading-path} that \(1-\tau=(g_H-c_H)/\gamma\), so
\[
J=\DT\,(1-\tau)=\frac{\DT\,(g_H-c_H)}{\gamma},\qquad \gamma=\bar\Gamma-\DT,
\]
where \(\bar\Gamma=\E_H\theta-\E_L\theta\) is the fixed class mean gap (Lemma~A.5). The \emph{balance condition} of
Appendix~A.4 is: the product \(\DT(r)\,[g_H(r)-c_H]\) is nondecreasing in \(r\).

The paper reads it in words: where class \(H\)'s gross profit exceeds the high cost, the spread must grow, in
proportion, at least as fast as the excess \(g_H-c_H\) shrinks. On a differentiable path, where \(g_H>c_H\):
\begin{align*}
\bigl(\DT\,(g_H-c_H)\bigr)'\ge0
&\iff \DT'\,(g_H-c_H)+\DT\,g_H'\ge0 \why{product rule; \(c_H\) fixed}\\
&\iff \frac{\DT'}{\DT}\ge\frac{-g_H'}{g_H-c_H} \why{divide by \(\DT(g_H-c_H)>0\); rearrange}
\end{align*}
The left side is the growth rate of the spread, and the right side is the rate at which the excess shrinks.

\paragraph{Part (c)(i).} Let the balance condition hold along a path. Then the strengths with \(J(r)\ge k\) form an
up-set, on which \(J\) is nondecreasing and rises strictly wherever \(\DT\) rises strictly. The strengths with
coexistence form an interval, at every prior and every slot.

\begin{derivation}[proof of part (c)(i)]
Let \(J(r)\ge k\) and \(r'>r\). Since \(\gamma(r)>0\) and \(J(r)\ge k>0\),
\(\DT(r)[g_H(r)-c_H]=\gamma(r)J(r)>0\). Then
\begin{align*}
J(r') &= \frac{\DT(r')\,[g_H(r')-c_H]}{\gamma(r')} \why{\(J=\DT(g_H-c_H)/\gamma\)}\\
&\ge \frac{\DT(r)\,[g_H(r)-c_H]}{\gamma(r')} \why{balance condition; \(\gamma(r')>0\)}\\
&\ge \frac{\DT(r)\,[g_H(r)-c_H]}{\gamma(r)} \why{\(0<\gamma(r')\le\gamma(r)\); positive numerator}\\
&= J(r)\ge k
\end{align*}
So \(\{J\ge k\}\) is an up-set, and \(J\) does not fall on it. If \(\DT(r')>\DT(r)\), then
\(\gamma(r')=\bar\Gamma-\DT(r')<\gamma(r)\), and the third line is strict.

By part (a), and because \(\tau\le M<1\) at every strength with coexistence, the condition \(k/(1-\tau)\le\DT\) is
\(J\ge k\). So the strengths with coexistence are
\[
\{\pi<\tau\}\cap\{J\ge k\}\cap\{\tau\le M\}\cap\{\DT\le\DT^{\rm N}\}:
\]
two up-sets and two down-sets, by Section~\ref{sec:coexistence-proof-b}. They form an interval. The balance condition
does not involve the prior or the slot, so the argument holds at every prior and every slot.
\end{derivation}

\paragraph{Part (c)(ii).} Let the interval of strengths \(I\) be closed and bounded, let \(\DT\) and \(\tau\) be
continuous in \(r\), and let \(\tau>\pi\) at the weakest incumbent of the path. Then the coexistence interval is
closed. When it is not empty, it starts at the weakest strength with \(J(r)\ge k\) and ends at the smaller of the
strongest strength with \(\tau(r)\le M\) and the strongest strength with \(\DT(r)\le\DT^{\rm N}\).

\emph{Proof.} The set \(\{\pi<\tau\}\) is an up-set that contains \(\min I\), so it is all of \(I\). Continuity makes
\(\{J\ge k\}\), \(\{\tau\le M\}\) and \(\{\DT\le\DT^{\rm N}\}\) closed subsets of \(I\). A closed up-set of \(I\) is
empty or runs from its least element to \(\max I\). A closed down-set is empty or runs from \(\min I\) to its largest
element. The intersection is therefore empty or the closed interval with the stated ends.

\paragraph{Part (c)(iii): uniform incumbents.} Let \(F_r=U[0,r]\) for \(r\) in a closed and bounded interval \(I\)
with \(p<\min I\). Suppose that every low-class value lies in \([p,\min I]\), every high-class value is at least
\(\max I\), and \(\E_H\theta\ge\max I+c_H\). Then \(\DT(r)[g_H(r)-c_H]\) is nondecreasing on \(I\), and \(J\) rises
strictly. For point classes at \(h\) and \(\ell\), the hypotheses read \(\ell\in[p,\min I]\), \(p<\min I\) and
\(\max I+c_H\le h\).

The proof has four steps. Fix \(r\in I\), and let \(R\) be uniform on \([0,r]\). Every \(r\in I\) has \(r>p\), so
\((r^2-p^2)/(2r^2)>0\).

\begin{derivation}[Step 1: the auction terms]
\emph{A high-class challenger.} Its value \(\theta\ge\max I\ge r\) beats every incumbent value, and it pays
\(\max\{p,R\}\). The mean payment adds up the reserve when \(R<p\) and \(R\) itself when \(R\ge p\):
\begin{align*}
\E\max\{p,R\} &= p\Pr(R<p)+\int_p^r y\,\frac{1}{r}\,dy \why{split at \(p\); density \(1/r\)}\\
&= \frac{p^2}{r}+\frac{r^2-p^2}{2r}=\frac r2+\frac{p^2}{2r} \why{\(\Pr(R<p)=p/r\); common denominator}
\end{align*}
\emph{A low-class challenger.} Its value \(\theta\in[p,\min I]\subseteq[p,r]\). It wins when \(R<\theta\) and keeps
\(\theta-\max\{p,R\}\). Its mean rent adds up these gains over the incumbent values below \(\theta\):
\begin{align*}
\E(\theta-\max\{p,R\})_+ &= (\theta-p)\frac{p}{r}+\int_p^\theta(\theta-y)\,\frac1r\,dy \why{\(R<p\): pays \(p\); \(p\le R<\theta\): pays \(R\)}\\
&= \frac{(\theta-p)\,p}{r}+\frac{(\theta-p)^2}{2r} \why{\(\int_p^\theta(\theta-y)\,dy=(\theta-p)^2/2\)}\\
&= \frac{(\theta-p)(2p+\theta-p)}{2r}=\frac{\theta^2-p^2}{2r} \why{factor \(\theta-p\)}
\end{align*}
The target receives \(\theta\) minus this rent: when the challenger wins it pays \(\max\{p,R\}\), and when it loses
the incumbent pays the challenger's bid \(\theta\), with zero rent. Averages over each class give
\begin{align*}
g_H &= \E_H\theta-\frac r2-\frac{p^2}{2r},\qquad g_L=\frac{\E_L\theta^2-p^2}{2r}, \why{the two lines above}\\
\DT &= t_H-t_L=\frac r2+\frac{p^2}{2r}-\E_L\theta+\frac{\E_L\theta^2-p^2}{2r} \why{\(t_H=\E\max\{p,R\}\), \(t_L=\E_L\theta-g_L\)}\\
&= \frac{r^2-2r\,\E_L\theta+\E_L\theta^2}{2r}=\frac{\E_L(r-\theta)^2}{2r} \why{\(p^2\) cancels; common denominator}
\end{align*}
So \(\gamma=g_H-g_L=\E_H\theta-r/2-\E_L\theta^2/(2r)\), and \(\DT+\gamma=\E_H\theta-\E_L\theta=\bar\Gamma\). The
derivatives are
\begin{align*}
\DT' &= \frac12-\frac{\E_L\theta^2}{2r^2}=\frac{r^2-\E_L\theta^2}{2r^2}\ge0, \why{\(\E_L\theta^2\le(\min I)^2\le r^2\)}\\
g_H' &= -\frac12+\frac{p^2}{2r^2}=-\frac{r^2-p^2}{2r^2},\qquad \gamma'=-\DT'. \why{\(\bar\Gamma\) is fixed}
\end{align*}
Also \(\DT(r)>0\) for every \(r>\min I\), because every low-class value is at most \(\min I<r\).
\end{derivation}

For point classes, \(\DT=(r-\ell)^2/(2r)\), the example's spread of Chapter~\ref{ch:example}.

\begin{derivation}[Step 2: three inequalities]
\emph{First, the excess \(g_H-c_H\) is large.} Since \(\E_H\theta-c_H\ge\max I\ge r\),
\begin{align*}
g_H-c_H &= \E_H\theta-c_H-\frac r2-\frac{p^2}{2r} \why{Step 1}\\
&\ge r-\frac r2-\frac{p^2}{2r}=\frac{r^2-p^2}{2r}=r\cdot\frac{r^2-p^2}{2r^2}>0 \why{\(\E_H\theta-c_H\ge r\)}
\end{align*}
\emph{Second, the spread grows at least in proportion to \(r\).}
\begin{align*}
r\DT'-\DT &= \frac r2-\frac{\E_L\theta^2}{2r}-\frac r2+\E_L\theta-\frac{\E_L\theta^2}{2r} \why{Step 1}\\
&= \E_L\theta-\frac{\E_L\theta^2}{r}=\E_L\Bigl[\theta\Bigl(1-\frac\theta r\Bigr)\Bigr]\ge0 \why{\(0\le\theta\le r\)}
\end{align*}
\emph{Third, the profit gap is positive.} With \(\E_H\theta\ge r+c_H\) and \(\E_L\theta^2\le r^2\),
\[
\gamma\ge(r+c_H)-\frac r2-\frac{r^2}{2r}=c_H>0 .
\]
Hence \(g_H>g_L\), and \(\bar\Gamma=\DT+\gamma>0\).
\end{derivation}

\begin{derivation}[Step 3: the product]
\begin{align*}
\bigl(\DT\,(g_H-c_H)\bigr)' &= \DT'\,(g_H-c_H)+\DT\,g_H' \why{product rule}\\
&= \DT'\,(g_H-c_H)-\frac{r^2-p^2}{2r^2}\,\DT \why{Step 1}\\
&\ge \DT'\,r\,\frac{r^2-p^2}{2r^2}-\frac{r^2-p^2}{2r^2}\,\DT \why{first inequality; \(\DT'\ge0\)}\\
&= \frac{r^2-p^2}{2r^2}\,\bigl(r\DT'-\DT\bigr)\ge0 \why{factor; second inequality}
\end{align*}
So \(\DT(r)[g_H(r)-c_H]\) is nondecreasing on \(I\).
\end{derivation}

\begin{derivation}[Step 4: the return \(J\)]
Lemma~A.5 gives \(J'\) on a differentiable path (Section~\ref{sec:trading-path}). Write \(c=(r^2-p^2)/(2r^2)>0\).
\begin{align*}
J' &= \frac{\bar\Gamma\,(g_H-c_H)\,\DT'+\DT\,\gamma\,g_H'}{\gamma^2} \why{Lemma~A.5}\\
&= \frac{\bar\Gamma\,(g_H-c_H)\,\DT'-c\,\DT\,\gamma}{\gamma^2} \why{\(g_H'=-c\)}\\
&\ge \frac{c\,\bigl(\bar\Gamma\,r\DT'-\DT\,\gamma\bigr)}{\gamma^2} \why{\(g_H-c_H\ge rc\); \(\bar\Gamma\DT'\ge0\)}\\
&\ge \frac{c\,\bigl(\bar\Gamma\,\DT-\DT\,\gamma\bigr)}{\gamma^2} \why{\(r\DT'\ge\DT\); \(\bar\Gamma>0\)}\\
&= \frac{c\,\DT\,(\bar\Gamma-\gamma)}{\gamma^2}=\frac{c\,\DT^2}{\gamma^2} \why{\(\bar\Gamma-\gamma=\DT\)}
\end{align*}
The bound is positive wherever \(\DT>0\), which includes every \(r>\min I\). So \(J\) rises strictly on \(I\).
\end{derivation}

\paragraph{What this says.} In this uniform family, a stronger incumbent raises the spread fast enough to beat the
fall in class \(H\)'s excess profit. Class \(H\) beats every incumbent, so its payment rises by less than one half
per unit of strength, while the spread rises at least in proportion to strength. Once a purchase that just reaches the
threshold covers \(k\), it covers \(k\) at every stronger incumbent of the family, and the coexistence set is one
interval at every prior. The example meets the hypotheses: its incumbent is uniform, \(\ell\) lies between \(p\) and
\(r_0\), and \(h\) exceeds \(r_2\) by more than \(c_H\) (Section~4.3 of the paper).

\begin{watchout}
Keep each hypothesis of part (c)(iii).
\begin{itemize}
\item \(\E_H\theta\ge\max I+c_H\) gives \(g_H-c_H\ge r(r^2-p^2)/(2r^2)\), the first inequality of Step~2. Without it the
excess can shrink faster than the spread grows, as in the instance of Lemma~A.5 where \(J\) falls.
\item Every low-class value in \([p,\min I]\) gives \(\DT'\ge0\) and \(r\DT'\ge\DT\).
\item Every high-class value at least \(\max I\) makes class \(H\) win for sure, so its payment is
\(r/2+p^2/(2r)\).
\end{itemize}
The two-interval path of Section~\ref{sec:coexistence-race} keeps class \(H\) winning for sure, with a value above
every incumbent value by more than \(c_H\), and still splits the coexistence set. The uniform shape matters.
\end{watchout}

% =====================================================================
\section{The race between the spread and the threshold}\label{sec:coexistence-race}

Section~\ref{sec:trading-race} described a race. Along a path, while \(g_L\le c_H\) and the profit gap stays
positive, a stronger incumbent never narrows the spread and never lowers the threshold. A wider spread raises what a
purchase that recruits earns. A higher threshold makes recruitment harder. In the coexistence region the race is the
number \(J=\DT(1-\tau)\): its first factor never falls and its second never rises.

So as the incumbent strengthens, the economy can enter the region (when \(J\) rises to \(k\)), leave it (when \(J\)
falls below \(k\)), and enter it again (when \(J\) rises to \(k\) once more). The other three conditions cannot
produce this: \(\{\pi<\tau\}\) can only switch on, and \(\{\tau\le M\}\) and \(\{\DT\le\DT^{\rm N}\}\) can only switch
off.

\paragraph{A path with two intervals.} Appendix~A.4 builds such a path. It keeps every primitive of the example except
the incumbent law, takes the prior \(1/2\) and a slot with \(\tau(r_1)<M(1/2)<\tau(r_2)\). The incumbent law is a mix
of two uniform laws of the same small width, a lower one and an upper one, with fixed weights.
\begin{itemize}
\item On the first half of the path only the lower law moves up, and its values stay between \(p\) and \(\ell\). The
incumbent strengthens only where it loses to both classes. Both gross profits fall, the spread stays constant, the
threshold rises, and \(J\) falls.
\item On the second half only the upper law moves up, with values between \(\ell\) and \(h\). The incumbent
strengthens between the class values, and the spread widens.
\end{itemize}
A quiet market and an informative market coexist exactly on an initial stretch and on a final stretch of the path,
with a gap between them. The two stretches end and start at the two strengths where \(J=k\). The strengths with a sure
purchase form one interval inside the second stretch, and on the gap the spread lies below \(\DT^{\rm F}\), as part
(b) requires.

\begin{derivation}[why the first half moves \(J\) down]
For point classes, \(D(y)=G_L(y)-G_H(y)=\1\{\ell\le y<h\}\), \(1-G_H(y)=\1\{y<h\}\) and \(1-G_L(y)=\1\{y<\ell\}\). On
the first half, \(F_r-F_{r'}\) is zero outside \([p,\ell)\) (Section~\ref{sec:toolkit-fosd}). Lemma~A.5 gives
\begin{align*}
\DT(r')-\DT(r) &= \int_p^\infty[F_r-F_{r'}]\,D\,dy=0 \why{\(D=0\) below \(\ell\)}\\
g_j(r)-g_j(r') &= \int_p^\infty[F_r-F_{r'}]\,[1-G_j]\,dy=\int_p^\ell[F_r-F_{r'}]\,dy \why{\(1-G_j=1\) below \(\ell\), both \(j\)}
\end{align*}
So \(g_H\) and \(g_L\) fall by the same positive amount, and \(\gamma=g_H-g_L\) stays fixed. Then
\(\tau=(c_H-g_L)/\gamma\) rises, because its numerator rises and its denominator stays fixed, and
\(J=\DT(1-\tau)\) falls, because \(\DT\) stays fixed.
\end{derivation}

The balance condition fails on the first half: the spread is constant while \(g_H-c_H\) falls, so the product
\(\DT(g_H-c_H)\) falls. This is what part (c) rules out. Under the balance condition, \(\{J\ge k\}\) is an up-set
(part (c)(i)). Once a purchase that just reaches the threshold covers \(k\), it covers \(k\) at every stronger
incumbent. The economy can then leave the region only through \(\tau\le M\) or \(\DT\le\DT^{\rm N}\), and it cannot
return through either. So it enters once and leaves once.

\begin{checkyourself}\label{q:coexistence-gap}
On a path, suppose that over a stretch of strengths the spread stays constant while \(g_H\) and \(g_L\) fall by the
same positive amount, with \(g_L\le c_H<g_H\). Show that \(\tau\) rises and \(J\) falls on the stretch. Show that the
balance condition fails there. Which condition of the region can switch off on this stretch and switch on again later?
\end{checkyourself}

% =====================================================================
\section{Corollary 2: coexistence in the example}\label{sec:coexistence-example}

The example has point classes at \(h\) and \(\ell\), a reserve \(p\), and an incumbent uniform on \([0,r]\). The path
is \(r\in[r_0,r_2]\), and on it \(p<\ell<r<h\). Chapter~\ref{ch:example} derives
\[
\DT(r)=\frac{(r-\ell)^2}{2r},\qquad \tau(r)=\frac{2rc_H-(\ell^2-p^2)}{2rh-r^2-\ell^2}.
\]
On the path \(\DT(r)>0\), and \(\gamma=(2rh-r^2-\ell^2)/(2r)>0\) because \(2rh-r^2-\ell^2-(r^2-\ell^2)=2r(h-r)>0\) and
\(r^2-\ell^2>0\). So every law of the path meets the hypotheses of Proposition~3. Take the prior \(1/2\) and a slot in
the band of Corollary~1: \(0<\lambda<1\), \(0<\beta<1/2\) and \(\tau(r_1)<M(1/2)<\tau(r_2)\). Write \(M=M(1/2)\) and
\(m=1-M\).

\begin{derivation}[the formulas at the prior \(1/2\)]
\begin{align*}
1-\mu_+(s) &= \frac{b/2}{b+s/2}=\frac{b}{2b+s} \why{\(\pi=1/2\); multiply above and below by \(2\)}\\
\mu_-(s) &= \frac{b/2}{b+s/2}=\frac{b}{2b+s} \why{the same steps}\\
1-M &= \frac{b}{2b+\lambda}=\frac{1}{2+u}=m \why{\(s=\lambda\); divide above and below by \(b\)}\\
s^* &= \frac{b\,(\tau-1/2)}{(1/2)(1-\tau)}=\frac{b\,(2\tau-1)}{1-\tau} \why{multiply above and below by \(2\)}
\end{align*}
\end{derivation}

At the prior \(1/2\) the doubt left after a purchase at rate \(s\) equals the doubt left after a sale at rate \(s\).
This symmetry makes the sale side and the purchase side easy to compare.

\paragraph{A root function.} The paper writes \(\fr(y)=\ell+y+\sqrt{y^2+2\ell y}\) for the root above \(\ell\) of
\(\DT(r)=y\), for \(y>0\):
\begin{align*}
\frac{(r-\ell)^2}{2r} &= y \why{\(\DT(r)=y\)}\\
r^2-2(\ell+y)\,r+\ell^2 &= 0 \why{multiply by \(2r\); expand; collect}\\
r &= \ell+y\pm\sqrt{(\ell+y)^2-\ell^2}=\ell+y\pm\sqrt{y^2+2\ell y} \why{quadratic formula}
\end{align*}
The root with the minus sign lies below \(\ell\), because \(\sqrt{y^2+2\ell y}>y\). Above \(\ell\) the spread rises
strictly (\(\DT'=\tfrac12-\ell^2/(2r^2)>0\)), so for \(r>\ell\), \(\DT(r)\le y\) holds exactly when \(r\le\fr(y)\).

\subsection{The marked strengths}

Each marked strength is the strength at which one order just earns the trading cost, or at which the threshold meets
the top posterior.
\begin{itemize}
\item \(r_A\): the root of \(J(r)=k\) in \((r_0,r_2)\). A purchase that just reaches the threshold and recruits earns
\(k\).
\item \(r_F=\fr\bigl(k/(1-M)\bigr)\): \(\DT(r)(1-M)=k\). A full purchase that recruits earns \(k\).
\item \(r_N=\fr(2k/\rho)\): \(\rho\DT(r)/2=k\). A unit order against an uninformative price earns \(k\): entry
\(\rho\) and the posterior \(1/2\) give \(\rho\DT\cdot\tfrac12\).
\item \(r_J\): the root of \(\rho J(r)=k\) in \((r_0,r_2)\). A purchase that just reaches the threshold but does not
recruit, at entry \(\rho\), earns \(k\).
\item \(r_S=\fr\bigl(k/(\rho m)\bigr)\): \(\rho\DT(r)m=k\). A full sale earns \(k\).
\item \(r_C\): the root of \(\tau(r)=M\) in \((r_1,r_2)\), that is of \(B_r(M)=c_H\). This is the entry ceiling:
above it even the top posterior does not recruit.
\item \(r_*\): the root of \(\rho\DT(r)\tau(r)=k\) in \((r_0,r_2)\). It matters only at other priors
(Section~\ref{sec:coexistence-priors}).
\end{itemize}
The strengths \(r_A\), \(r_N\), \(r_J\) and \(r_*\) do not depend on the slot. The strengths \(r_F\), \(r_S\) and \(r_C\)
depend on it through \(M\). Chapter~\ref{app:numbers} lists the values at the example's slot.

\begin{paperresult}{Lemma A.10 (The example path)}
On the path \([r_0,r_2]\) of the example:
\begin{enumerate}[label=(\alph*),leftmargin=2em]
\item \(\tau\) is strictly increasing, with \(\tau(r_0)>1/2\) and \(\tau(r_2)<1-\rho\).
\item The path is a uniform incumbent of Proposition~3(c)(iii), with \(h-r_2-c_H>0\). The product
\(\DT(r)\,[g_H(r)-c_H]\) has a positive derivative at every \(r\). Hence \(J\) is strictly increasing, with
\(J(r_0)<k\) and \(\rho J(r_2)>k\).
\item The product \(\DT(r)\,[g_H(r)-c_H]^2\) has a positive derivative at every \(r\), and \(\DT(r)\,[1-\tau(r)]^2\) is
strictly increasing on the whole path.
\item \(\DT(r)\,\tau(r)\) is strictly increasing.
\end{enumerate}
\end{paperresult}

The paper states each value and each margin in this lemma as an exact number. The guide keeps only the inequalities
that the proofs use; Chapter~\ref{app:numbers} lists the numbers.

\begin{derivation}[proof of Lemma A.10]
\emph{(a)} Write \(\tau=N/D\) with \(N=2rc_H-\ell^2+p^2\) and \(D=2rh-r^2-\ell^2>0\). The derivative of \(\tau\) has
the sign of \(N'D-ND'\):
\begin{align*}
N'D-ND' &= 2c_H\,(2rh-r^2-\ell^2)-\bigl(2rc_H-\ell^2+p^2\bigr)(2h-2r) \why{\(N'=2c_H\), \(D'=2h-2r\)}\\
&= 4c_Hrh-2c_Hr^2-2c_H\ell^2-4c_Hrh+4c_Hr^2+2(\ell^2-p^2)(h-r) \why{expand}\\
&= 2c_H\,(r^2-\ell^2)+2(\ell^2-p^2)(h-r)>0 \why{collect; \(\ell<r<h\), \(\ell\ge p\)}
\end{align*}
The values of the example give \(\tau(r_0)>1/2\) and \(\tau(r_2)<1-\rho\).

\emph{(b)} The interval \(I=[r_0,r_2]\) has \(p<\ell<r_0\) and \(r_2+c_H<h\), so part (c)(iii) applies:
\(\DT(g_H-c_H)\) is nondecreasing and \(J\) rises strictly. The paper also bounds the derivative directly. From
\(\DT'=(r-\ell)(r+\ell)/(2r^2)\), the ratio \(\DT/\DT'\) is \(r(r-\ell)/(r+\ell)\), so
\begin{align*}
\bigl(\DT(g_H-c_H)\bigr)' &= \DT'(g_H-c_H)+\DT\,g_H' \why{product rule}\\
&= \DT'\Bigl[g_H-c_H+\frac{\DT}{\DT'}\,g_H'\Bigr] \why{factor \(\DT'>0\)}\\
&= \DT'\Bigl[g_H-c_H-\frac{(r^2-p^2)(r-\ell)}{2r(r+\ell)}\Bigr] \why{\(g_H'=-(r^2-p^2)/(2r^2)\)}
\end{align*}
In the bracket, \(g_H-c_H\) falls strictly and the subtracted term rises strictly, as the product of the positive
increasing functions \((r^2-p^2)/(2r)\) and \((r-\ell)/(r+\ell)\). So the bracket is smallest at \(r_2\), and the
example's values make it positive there. The endpoint values \(J(r_0)<k\) and \(\rho J(r_2)>k\) follow from the
closed form of \(J\) at the example's values.

\emph{(c)} In the same way,
\[
\bigl(\DT(g_H-c_H)^2\bigr)'=(g_H-c_H)\,\DT'\Bigl[g_H-c_H-\frac{(r^2-p^2)(r-\ell)}{r(r+\ell)}\Bigr],
\]
and the bracket, strictly decreasing for the same reasons, is positive at \(r_2\) at the example's values. Since
\(g_H-c_H>0\), the derivative is positive. Then \(\DT(1-\tau)^2=\DT(g_H-c_H)^2/\gamma^2\) has a positive, strictly
increasing numerator and a positive, strictly decreasing denominator (\(\gamma'=-\DT'<0\)), so it rises strictly.

\emph{(d)} \(\DT\) and \(\tau\) are positive and strictly increasing on the path, by \(\DT'>0\) and (a), so their
product is strictly increasing.
\end{derivation}

By part (b), \(J(r_0)<k<\rho J(r_2)<J(r_2)\) and \(\rho J(r_0)<J(r_0)<k\). So the roots \(r_A\) of \(J=k\) and \(r_J\) of
\(\rho J=k\) exist in \((r_0,r_2)\), and they are unique because \(J\) rises strictly. The root \(r_C\) is unique
because \(\tau\) rises strictly, and \(r_*\) because \(\DT\tau\) does.

\begin{paperresult}{Lemma A.11 (The order of the marked strengths)}
For every slot with \(0<\lambda<1\), \(0<\beta<1/2\) and \(\tau(r_1)<M(1/2)<\tau(r_2)\),
\[
r_0<r_A<r_F<r_N<r_J<r_S<r_1<r_C<r_2 .
\]
\end{paperresult}

\begin{derivation}[proof of Lemma A.11]
Every strength in the chain exceeds \(\ell\), where \(\DT\) rises strictly. Besides Lemma~A.10 and the band, the proof
uses two facts of the example's values: \(\rho J(r_1)>k\) and \(\rho^2\DT(r_1)>k\).
\begin{itemize}
\item \(r_0<r_A<r_1\): \(J(r_0)<k<\rho J(r_1)<J(r_1)\), and \(J\) rises strictly.
\item \(r_A<r_F\): \(\DT(r_A)<k/(1-M)\) reads \((1-M)\DT(r_A)<k=(1-\tau(r_A))\DT(r_A)\), that is \(M>\tau(r_A)\). This
holds because \(M>\tau(r_1)>\tau(r_A)\), since \(r_A<r_1\).
\item \(r_F<r_N\): \(k/(1-M)<2k/\rho\) reads \(M<1-\rho/2\). This holds because \(M<\tau(r_2)<1-\rho\), by
Lemma~A.10(a).
\item \(r_N<r_J\): \(\rho\DT(r_J)[1-\tau(r_J)]=k\) and \(\tau(r_J)>\tau(r_0)>1/2\) give
\(\rho\DT(r_J)/2>\rho\DT(r_J)[1-\tau(r_J)]=k\), that is \(\DT(r_J)>2k/\rho=\DT(r_N)\).
\item \(r_J<r_S\): \(\DT(r_J)<k/(\rho m)\) reads \(\rho(1-M)\DT(r_J)<k=\rho(1-\tau(r_J))\DT(r_J)\), that is
\(M>\tau(r_J)\). The fact \(\rho J(r_1)>k\) gives \(r_J<r_1\), so \(M>\tau(r_1)>\tau(r_J)\).
\item \(r_S<r_1\): the claim is \(\rho m\DT(r_1)>k\). By Lemma~A.10(a), \(m=1-M>1-\tau(r_2)>\rho\), so
\(\rho m\DT(r_1)>\rho^2\DT(r_1)>k\).
\item \(r_1<r_C<r_2\): this is the band, since \(\tau\) rises strictly.
\end{itemize}
\end{derivation}

Read the chain as a walk along the path. First a purchase at the threshold that recruits starts to pay (\(r_A\)),
then a full purchase that recruits (\(r_F\)), then an order against a silent price (\(r_N\)), then a purchase at the
threshold without recruitment (\(r_J\)), then a full sale (\(r_S\)). Far to the right the threshold passes the top
posterior (\(r_C\)).

\subsection{The equilibrium set}

At the prior \(1/2\), the low class sells at the rate of Lemma~A.8, which at this prior reads
\[
\hat s_-(r)=0\ \text{on}\ [r_0,r_N],\qquad \hat s_-(r)=b\,[\rho\DT(r)/k-2]\ \text{on}\ (r_N,r_S),\qquad
\hat s_-(r)=\lambda\ \text{on}\ [r_S,r_2].
\]
The middle formula is \(b(\pi\rho\DT/k-1)/(1-\pi)\) at \(\pi=1/2\). Class \(H\)'s pair takes one of five forms:
\begin{itemize}
\item \textbf{N}, a quiet market: \(s_+=0\) and \(e(1)=\rho\);
\item \textbf{R}, a mixed purchase that recruits: \(s_+=s_R(r):=b[\DT(r)/k-2]\) and \(e(1)=1\);
\item \textbf{F}, a full purchase: \(s_+=\lambda\), with \(e(1)=1\) for \(r\le r_C\) (at \(r_C\) because an
indifferent high-cost challenger prepares) and \(e(1)=\rho\) for \(r>r_C\);
\item \textbf{Q}, mixed orders at the entry floor: \(s_+=s_Q(r):=b[\rho\DT(r)/k-2]=\hat s_-(r)\) and \(e(1)=\rho\);
\item \textbf{T}, the threshold pair: \(s_+=s^*(r)\) and \(e(1)=k/J(r)\), which under the default rule requires
\(J(r)=k\), that is \(r=r_A\).
\end{itemize}
The forms R, Q and T need an indifferent investor who may mix. In Q both classes trade at the same rate, no order
recruits, and entry stays at the floor.

\begin{paperresult}{Corollary 2 (Coexistence in the example, the full statement of Appendix A.4)}
Take the example at the prior \(1/2\), a slot with \(0<\lambda<1\), \(0<\beta<1/2\) and \(\tau(r_1)<M(1/2)<\tau(r_2)\),
and the default tie rules. The equilibrium set at each strength of the path is
\begin{center}\upshape
\begin{tabular}{@{}ll@{}}
\toprule
Strength & Equilibria\\
\midrule
\(r_0\le r<r_A\) & \textbf{N}\\
\(r=r_A\) & \textbf{N}; \textbf{T} with entry one after a purchase\\
\(r_A<r<r_F\) & \textbf{N}; \textbf{R}\\
\(r_F\le r\le r_N\) & \textbf{N}; \textbf{F}\\
\(r_N<r<r_J\) & \textbf{F}; \textbf{Q}\\
\(r_J\le r\le r_2\) & \textbf{F}\\
\bottomrule
\end{tabular}
\end{center}
In each, the low class sells at the rate \(\hat s_-(r)\). Entry is \(\rho\) in a quiet market and in mixed orders at
the entry floor; \(\Ent_{\rm R}(r)=\rho+(1-\rho)\,b\DT(r)/(2k)\) in a mixed purchase that recruits, and
\(\Ent_{\rm R}(r_A)\) in the threshold pair at \(r_A\); and \(\Ent_{\rm F}=\rho+(1-\rho)(b+\lambda/2)\) in a full
purchase up to \(r_C\), with \(\Ent_{\rm F}=\rho\) above \(r_C\).
\begin{enumerate}[label=(\alph*),leftmargin=2em]
\item A quiet market and an informative market coexist exactly on \([r_A,r_N]\).
\item Along the informative branch, entry is \(\rho+(1-\rho)\min\{b\DT(r)/(2k),\,b+\lambda/2\}\) on \([r_A,r_C]\) and
\(\rho\) on \((r_C,r_2]\). It is continuous on \([r_A,r_C]\), rises strictly on \([r_A,r_F]\), is constant on
\([r_F,r_C]\), and returns to the floor \(\rho\) above the entry ceiling \(r_C\).
\item The equilibrium set is nonempty at every strength.
\end{enumerate}
\end{paperresult}

The main text states a short version: an equilibrium exists at every strength from \(r_0\) to \(r_2\), coexistence
holds exactly on \([r_A,r_N]\), and in the informative market that coexists with a quiet market entry is
\(\rho+(1-\rho)\min\{b\DT(r)/(2k),\,b+\lambda/2\}\), rising strictly on \([r_A,r_F]\) and constant on \([r_F,r_N]\);
with a full purchase, entry keeps this value up to \(r_C\) and equals \(\rho\) above it. The three economies of
Proposition~2 sit at \(r_0\), \(r_1\) and \(r_2\), and the table fills in the strengths between them: one quiet market
at \(r_0\), one full purchase that recruits at \(r_1\), and one full purchase at floor entry at \(r_2\).

\paragraph{The proof.} Appendix~A.4 proves the corollary in eight steps. The guide groups them in four.

\begin{derivation}[Step 1: the frame]
By Lemma~A.10(a), \(\tau(r)\ge\tau(r_0)>1/2=\pi\) on the whole path, so Lemmas~A.7 to~A.9 apply at every strength.
Every equilibrium combines a pair of Lemma~A.7 with a rate of Lemma~A.8. Recall \(1-M=1/(2+u)\) and \(ub=\lambda\).
\end{derivation}

\begin{derivation}[Step 2: the high-class pairs along the path]
\emph{(N)} holds exactly when \(\rho\DT/2\le k\), that is \(r\le r_N\).

\emph{(R).} The equation \(\DT\,b/(2b+s)=k\) gives \(2b+s=b\DT/k\), so \(s=s_R(r)\). Then
\begin{align*}
\mu_+(s_R)>\tau &\iff \frac{k}{\DT}<1-\tau \why{\(1-\mu_+(s_R)=b/(2b+s_R)=k/\DT\)}\\
&\iff k<J(r)\iff r>r_A \why{multiply by \(\DT>0\); \(J\) rises strictly}\\
s_R<\lambda &\iff \frac{\DT}{k}<2+u \why{divide \(b(\DT/k-2)<\lambda\) by \(b\); add \(2\)}\\
&\iff \DT(1-M)<k\iff r<r_F \why{\(1-M=1/(2+u)\)}
\end{align*}
The rate is positive for \(r>r_A\): there \(J>k\) and \(\tau>1/2\) give \(\DT/2>\DT(1-\tau)>k\), so \(\DT/k>2\). This
check is the guide's own. So (R) holds exactly on \((r_A,r_F)\).

\emph{(Q).} The equation \(\rho\DT\,b/(2b+s)=k\) gives \(s=s_Q(r)\). Then \(s_Q>0\) exactly when \(\rho\DT>2k\), that is
\(r>r_N\); \(\mu_+(s_Q)<\tau\) exactly when \(k/(\rho\DT)>1-\tau\), that is \(\rho J(r)<k\), that is \(r<r_J\); and
\(s_Q<\lambda\) exactly when \(\rho\DT(1-M)<k\), that is \(r<r_S\). Since \(r_J<r_S\), (Q) holds exactly on
\((r_N,r_J)\).

\emph{(T).} \(s^*\in(0,\lambda)\) exactly when \(1/2<\tau(r)<M\), that is \(r<r_C\). Also \(\rho\le k/J\le1\) exactly when
\(\rho J\le k\le J\), that is \(r_A\le r\le r_J\), and \(k/J=1\) exactly when \(r=r_A\). Since \(r_J<r_C\), (T) holds at
\(r_A\) under the default challenger rule.

\emph{(F).} For \(r<r_C\), \(M>\tau(r)\), so \(e(1)=1\), and the condition is \(\DT(1-M)\ge k\), that is \(r\ge r_F\). For
\(r>r_C\), \(M<\tau(r)\), so \(e(1)=\rho\), and \(\rho\DT(1-M)\ge k\) holds because \(r>r_C>r_S\). At \(r=r_C\), the
indifferent challenger prepares, \(e(1)=1\), and the condition holds because \(\DT(r_C)(1-M)>\rho\DT(r_C)(1-M)>k\). So
(F) holds exactly when \(r\ge r_F\).
\end{derivation}

\begin{derivation}[Step 3: the low-class rate and the rule \(s_+\ge s_-\)]
At the prior \(1/2\), Lemma~A.8 gives \(s_-=0\) exactly when \(\rho\DT/2\le k\), that is \(r\le r_N\); \(s_-=\lambda\)
exactly when \(\rho\DT m\ge k\), that is \(r\ge r_S\); and \(\hat s_-(r)\) in between.

For \(r\le r_N\), \(s_-=0\le s_+\). For \(r>r_N\), the pairs of Step~2 under the default rules are (F), with
\(s_+=\lambda\ge s_-\), and (Q), with \(s_+=s_Q(r)=\hat s_-(r)\); the pairs (R) and (T) occur only below \(r_F<r_N\). So
\(s_+\ge s_-\) in every combination, and Lemma~A.9 makes the equilibrium set the set of all these combinations, with
entry \(\rho\) after the zero order and a sale.
\end{derivation}

\begin{derivation}[Step 4: the set, entry and existence]
\emph{The set.} Put Step~2 together with the order of Lemma~A.11. On \([r_0,r_A)\) only (N) holds. At \(r_A\), (N) and
(T). On \((r_A,r_F)\), (N) and (R). On \([r_F,r_N]\), (N) and (F). On \((r_N,r_J)\), (F) and (Q). On \([r_J,r_2]\), only
(F). This is the table, and part (a) follows. On \([r_0,r_N]\) it is Proposition~3(a) on the path, since there
\(\DT\le\DT^{\rm N}\), \(1/2<\tau<M\), and \(J\ge k\) exactly when \(r\ge r_A\).

\emph{Entry.} With \(e(0)=e(-1)=\rho\) and \(\Pr(X=1)=\tfrac12(b+s_+)+\tfrac12b=b+s_+/2\),
\begin{align*}
\Ent &= \rho\,[1-\Pr(X=1)]+e(1)\Pr(X=1)=\rho+[e(1)-\rho]\,(b+s_+/2) \why{total probability over orders}\\
b+\frac{s_R}{2} &= b+\frac b2\Bigl(\frac{\DT}{k}-2\Bigr)=\frac{b\,\DT}{2k} \why{form R}\\
\frac{b\,\DT(r_F)}{2k} &= \frac{b\,(2+u)}{2}=b+\frac{\lambda}{2} \why{\(\DT(r_F)=k(2+u)\); \(ub=\lambda\)}
\end{align*}
So a mixed purchase that recruits has \(\Ent_{\rm R}(r)=\rho+(1-\rho)b\DT(r)/(2k)\), which rises strictly because \(\DT\)
does, and it meets the full purchase at \(r_F\) with no jump. A full purchase has \(b+\lambda/2\), so
\(\Ent_{\rm F}=\rho+(1-\rho)(b+\lambda/2)\) with \(e(1)=1\) up to \(r_C\), and \(\Ent_{\rm F}=\rho\) above \(r_C\). The forms
N and Q have \(e(1)=\rho\), so entry is \(\rho\). At \(r_A\), \(J=k\) gives \(\DT/k=1/(1-\tau)\), so
\(s_R=b[1/(1-\tau)-2]=b(2\tau-1)/(1-\tau)=s^*\): the threshold pair is the limit of the mixed purchase, with entry
\(\Ent_{\rm R}(r_A)\). This gives part (b).

\emph{Existence.} A quiet market exists on \([r_0,r_N]\) and a full purchase on \([r_F,r_2]\). Since \(r_F<r_N\), every
strength has an equilibrium. This is part (c). Proposition~3(d) gives existence at the prior \(1/2\) for a general
incumbent law with \(\tau>1/2\).
\end{derivation}

\begin{intuition}[a walk along the path]
At the weak end the spread is so narrow that even a purchase that just reaches the threshold, and recruits, does not
cover \(k\): only the quiet market exists. At \(r_A\) that purchase covers \(k\) exactly, and the informative market
appears as the threshold pair. On \((r_A,r_F)\) class \(H\) mixes, and its purchase rate and entry rise with the spread.
From \(r_F\) class \(H\) buys for sure, and entry stays on a plateau. At \(r_N\) an order against a silent price starts
to pay, so the quiet market ends and coexistence ends with it. Beyond \(r_N\) the low class starts to sell, and from
\(r_S\) it sells for sure. At \(r_C\) the threshold passes the top posterior, and the full purchase stops recruiting.
\end{intuition}

% =====================================================================
\section{Reading the figure}\label{sec:coexistence-figure}

\begin{figure}[htbp]
\centering
\includegraphics[width=\linewidth]{coexistence}
\caption{Entry and the coexistence region along the example's path, at the prior \(1/2\), under the default tie
rules. Panels (a) and (c) plot entry against the strength \(r\) in each equilibrium of Corollary~2; panel (c)
enlarges the strengths around the shaded coexistence interval \([r_A,r_N]\). Panel (b) draws the region of
Proposition~3(a) in the plane of \(\DT\) and \(\tau\), with the curve \((\DT(r),\tau(r))\) of the path. The labels N,
R, F, Q and T name the quiet market, the mixed purchase that recruits, the full purchase, mixed orders at the entry
floor and the threshold pair. Chapter~\ref{app:numbers} lists the values of the marked strengths.}
\label{fig:coexistence-path}
\end{figure}

The paper's notes give one rule for the markers: a filled marker belongs to its curve, and an open marker does not.

\paragraph{Panel (a): the whole path.} The quiet market N is a flat line at the floor \(\rho\), from \(r_0\) to \(r_N\).
The informative branch starts at \(r_A\) above the floor, rises along the mixed purchase R up to \(r_F\), and then stays
on the plateau of the full purchase F up to the entry ceiling \(r_C\). At \(r_C\) the filled marker on the plateau shows
that the full purchase still recruits there, because an indifferent high-cost challenger prepares. Above \(r_C\) the
full purchase has entry \(\rho\): the open marker at the floor. The dashed curve that falls from the start of R is the
threshold pair T under the other challenger rule (Section~\ref{sec:coexistence-ties}). So entry along the informative
branch first rises with the spread, then stays flat, then drops to the floor when even the top posterior stops
recruiting.

\paragraph{Panel (c): around the coexistence interval.} The shaded band runs from \(r_A\) to \(r_N\). Inside it two
equilibria exist at every strength: N at the floor, and R or F above it. The step from R to F at \(r_F\) has no jump.
To the right of \(r_N\), the quiet market is gone, and the mixed orders at the entry floor Q run at the floor from
\(r_N\) to \(r_J\), next to the full purchase. The labels above the plateau show what the low class does in the full
purchase: it abstains up to \(r_N\), mixes from \(r_N\) to \(r_S\) (``mixed sales''), and sells for sure from \(r_S\)
(``full trade''). The dashed curve T falls from the level of R at \(r_A\) to the floor at \(r_J\).

\paragraph{Panel (b): the region and the path.} The axes are the spread \(\DT\) and the threshold \(\tau\). The region
is the set with \(1/2<\tau\le M(1/2)\) and \(k/(1-\tau)\le\DT\le2k/\rho\), at the example's slot: the shape of
Figure~\ref{fig:coexistence-region}. Its edges:
\begin{itemize}
\item the dotted lower edge \(\tau=1/2\) lies outside the region: at \(\tau=\pi\), the quiet market under the default
challenger rule has entry one after every order, the case of Step~1 in Section~\ref{sec:coexistence-proof-a};
\item the top edge \(\tau=M(1/2)\) belongs to it: there a full purchase reaches the threshold, and the indifferent
challenger prepares;
\item the left edge, the curve \(\DT=k/(1-\tau)\), belongs to it: there the threshold pair recruits, again because the
indifferent challenger prepares;
\item the right edge \(\DT=2k/\rho\) belongs to it: there the investor in the quiet market is indifferent, and it may
abstain.
\end{itemize}
The dashed line \(\DT=k/(1-M(1/2))\) divides the mixed purchase R, on its left, from the full purchase F. The curve
traces \((\DT(r),\tau(r))\) from \(r_0\). Both coordinates rise with \(r\), and the threshold rises slowly. The curve
enters the region through the left edge at \(r_A\), crosses the dashed line at \(r_F\), leaves the region through the
right edge at \(r_N\), and continues toward \(r_1\) and \(r_2\) beyond the panel. It never touches the top edge on this
stretch, because \(r_C\) lies far to the right of \(r_N\).

% =====================================================================
\section{Other tie rules}\label{sec:coexistence-ties}

Proposition~3 and Corollary~2 take the default tie rules. Appendix~A.4 states what the other rules change. The
reasons all come from Lemmas~A.7 to~A.9, which hold under each pair of rules.

\paragraph{A general incumbent law.}
\begin{itemize}
\item \emph{The investor may mix, and the challenger may prepare with any probability.} Coexistence holds exactly in
the region of part (a), or when \(\tau=\pi\) and \(k/(1-\pi)<\DT\le\DT^{\rm N}\). At \(\tau=\pi\), entry \(\rho\) after
every order is allowed at the belief \(\pi=\tau\), which gives a quiet market; and a purchase at a positive rate has a
posterior above \(\tau\), so it recruits, and it pays when \(\DT(1-\pi)>k\). Inside the region, wherever \(\tau<M\) and
\(J>k\), the threshold pair with \(e(1)=k/J\) is a further informative market. Its entry lies strictly between \(\rho\)
and \eqref{eq:coexistence-entry}: \(e^*=k/J<1\), and \(e^*>\rho\) because \(\rho J<\rho\DT(1-\pi)\le k\). Where
\(\tau=M\), a full purchase is an equilibrium with each entry after a purchase in \([k/J,1]\).
\item \emph{The investor abstains, and the challenger prepares.} Coexistence holds exactly when \(\pi<\tau\le M\) and
\(\DT^{\rm F}<\DT\le\DT^{\rm N}\), and the informative market is the full purchase. Without mixes the pairs R, T and Q
drop out, and the full purchase needs the strict return \(\DT(1-M)>k\).
\item \emph{The investor abstains, and the challenger may prepare with any probability.} The laws with \(\tau=\pi\) and
\(\DT^{\rm F}<\DT\le\DT^{\rm N}\) join this set.
\end{itemize}

\paragraph{Proposition 3(c)(iv).} If, beyond the balance condition, \(\DT(r)[g_H(r)-c_H]^2\) is nondecreasing, then
\(\DT(r)[1-\tau(r)]^2\) is nondecreasing. When an indifferent investor may mix and an indifferent high-cost challenger
may prepare with any probability, the threshold pair is a pair of Lemma~A.7 exactly at the strengths with
\(\pi<\tau(r)<M\) and \(\rho J(r)\le k\le J(r)\). There its entry, with entry \(\rho\) after a sale and after the zero
order, is
\[
\Ent_{\rm T}(r)=\rho+(1-\pi)\,b\,\frac{k-\rho J(r)}{\DT(r)\,[1-\tau(r)]^2},
\]
and it is nonincreasing. Where \(\tau(r)=M\), the threshold rate is \(s^*=\lambda\), and the purchase is the full
purchase (F).

\begin{derivation}[the entry of the threshold pair, and why it falls]
\emph{The first claim.} \(\DT(1-\tau)^2=\DT(g_H-c_H)^2/\gamma^2\). The numerator is nonnegative and nondecreasing by
hypothesis, and \(1/\gamma^2\) is positive and nondecreasing because \(\gamma>0\) never rises. A product of two
nonnegative nondecreasing functions is nondecreasing.

\emph{The entry.} Entry is \(\rho+(e^*-\rho)(b+\pi s^*)\), by the entry derivation of
Section~\ref{sec:coexistence-informative} with \(e(1)=e^*\). Then
\begin{align*}
b+\pi s^* &= b+\frac{b\,(\tau-\pi)}{1-\tau}=\frac{b\,(1-\tau+\tau-\pi)}{1-\tau}=\frac{(1-\pi)\,b}{1-\tau} \why{\(s^*\); common denominator}\\
e^*-\rho &= \frac{k}{J}-\rho=\frac{k-\rho J}{\DT(1-\tau)} \why{\(e^*=k/J\), \(J=\DT(1-\tau)\)}\\
(e^*-\rho)(b+\pi s^*) &= (1-\pi)\,b\,\frac{k-\rho J}{\DT(1-\tau)^2} \why{multiply}
\end{align*}

\emph{It falls.} Take \(r<r'\) in this set. Both have \(J\ge k\), so part (c)(i) gives \(J(r')\ge J(r)\) and
\(0\le k-\rho J(r')\le k-\rho J(r)\). The first claim gives \(\DT(r')[1-\tau(r')]^2\ge\DT(r)[1-\tau(r)]^2>0\). A smaller
numerator over a larger denominator gives \(\Ent_{\rm T}(r')\le\Ent_{\rm T}(r)\).
\end{derivation}

In words: as the incumbent strengthens, \(J\) rises, so a smaller entry \(e^*=k/J\) after a purchase already pays the
investor. The threshold pair then recruits less and less, until at \(\rho J=k\) it recruits nobody beyond the floor.

\paragraph{The example at the prior \(1/2\).} Table~\ref{tab:coexistence-ties} copies the paper's Table~A.3.

\begin{table}[htbp]
\centering\small
\setlength{\tabcolsep}{4pt}
\begin{tabular}{@{}lllll@{}}
\toprule
 & \multicolumn{2}{l}{Investor may mix} & \multicolumn{2}{l}{Investor abstains}\\
\cmidrule(lr){2-3}\cmidrule(lr){4-5}
Result & prepares & any probability & prepares & any probability\\
\midrule
Proposition~2 & holds & holds & holds & holds\\
Quiet market \textbf{N} & \([r_0,r_N]\) & \([r_0,r_N]\) & \([r_0,r_N]\) & \([r_0,r_N]\)\\
Mixed purchase \textbf{R} & \((r_A,r_F)\) & \((r_A,r_F)\) & none & none\\
Threshold pair \textbf{T} & \(r_A\) only & \([r_A,r_J]\) & none & none\\
Full purchase \textbf{F} & \([r_F,r_2]\) & \([r_F,r_2]\) & \((r_F,r_N]\cup(r_S,r_2]\) & \((r_F,r_N]\cup(r_S,r_2]\)\\
Set on \((r_N,r_J)\) & \textbf{F}; \textbf{Q} & \textbf{F}; \textbf{Q}; \textbf{T} & empty & empty\\
Set at \(r_J\) & \textbf{F} & \textbf{F}; \textbf{T} & empty & empty\\
Set on \((r_J,r_S]\) & \textbf{F} & \textbf{F} & empty & empty\\
Entry after a purchase at \(r_C\) & \(1\) & any in \([\rho,1]\) & \(1\) & any in \([\rho,1]\)\\
Coexistence & \([r_A,r_N]\) & \([r_A,r_N]\) & \((r_F,r_N]\) & \((r_F,r_N]\)\\
\bottomrule
\end{tabular}
\caption{Which conclusions depend on the tie rules in the example, at the prior \(1/2\). The column heads give what an
indifferent investor does (may mix over its best orders, or abstains) and what an indifferent high-cost challenger
does (prepares, or may prepare with any probability). Source: Table~A.3 of the paper.}
\label{tab:coexistence-ties}
\end{table}

What changes, in words:
\begin{itemize}
\item Proposition~2 and the stretch of the quiet market do not depend on the tie rules.
\item The challenger rule matters only through the threshold pair and the point \(r_C\). When an indifferent
high-cost challenger may prepare with any probability, the threshold pair is also an equilibrium on \((r_A,r_J]\). Its
entry falls strictly from \(\Ent_{\rm R}(r_A)\) at \(r_A\) to \(\rho\) at \(r_J\), by Lemma~A.10(b) and~(c) and
Proposition~3(c)(iv). At \(r_C\), entry after a purchase in the full purchase can take any value in \([\rho,1]\). The
coexistence interval stays \([r_A,r_N]\).
\item The investor rule matters more. When an indifferent investor abstains, the mixed purchase and the threshold pair
disappear, and coexistence shrinks to \((r_F,r_N]\). Entry in the informative market then has no strict rise: it is
on the plateau wherever the market exists.
\item When an indifferent investor abstains, the equilibrium set is empty on \((r_N,r_S]\), under either challenger
rule. There the low class has no qualifying sale rate (Lemma~A.8). So an equilibrium at every strength of the path
needs the rule that an indifferent investor may mix.
\end{itemize}

At a general prior, when an indifferent investor may mix and an indifferent high-cost challenger may prepare with
any probability, the coexistence set of Section~\ref{sec:coexistence-priors} gains the single strength \(r_\pi\), where
\(\tau(r_\pi)=\pi\), for priors in \((\tau_A,\pi_{\max}]\). When an indifferent investor abstains, coexistence keeps only
the strengths at which a full purchase pays strictly.

\begin{checkyourself}\label{q:coexistence-abstain}
Let an indifferent investor abstain. Explain, step by step, why the example has no equilibrium at a strength \(r\) in
\((r_N,r_S]\). Which case of Lemma~A.8 applies? Why does the full purchase of class \(H\) not rescue the tuple?
\end{checkyourself}

% =====================================================================
\section{Coexistence at other priors}\label{sec:coexistence-priors}

Corollary~2 fixes the prior at \(1/2\). Appendix~A.4 also describes the coexistence set at every prior, in the
example, under the default tie rules (Corollary~A.2). This section gives its shape and the reason for each case.

Write \(r_N(\pi)=\fr\bigl(k/(\rho\max\{\pi,1-\pi\})\bigr)\), the strength at which an order against a silent price
starts to pay. Write \(r_C(\pi)\) for the largest strength with \(\tau(r)\le M(\pi)\), when one exists, and \(r_\pi\)
for the strength with \(\tau(r_\pi)=\pi\), when \(\tau(r_0)\le\pi\le\tau(r_2)\). Let \(\tau_A=\tau(r_A)\). Three priors
divide the cases:
\begin{itemize}
\item \(\pi_A^{\rm lo}\), with \(M(\pi_A^{\rm lo})=\tau_A\): the prior whose top posterior just reaches the threshold
at \(r_A\);
\item \(\pi_A^{\rm hi}=\tau_A\): the prior that equals the threshold at \(r_A\);
\item \(\pi_{\max}=\tau(r_*)\): the prior that equals the threshold at \(r_*\), where \(\rho\DT\tau=k\).
\end{itemize}
The odds form gives \(\pi_A^{\rm lo}\): its odds times \(1+u\) equal the odds of \(\tau_A\), so
\begin{align*}
\frac{\pi_A^{\rm lo}}{1-\pi_A^{\rm lo}} &= \frac{\tau_A}{(1-\tau_A)(1+u)} \why{divide the odds of \(\tau_A\) by \(1+u\)}\\
\pi_A^{\rm lo} &= \frac{\tau_A}{\tau_A+(1-\tau_A)(1+u)} \why{probability from odds}
\end{align*}
The priors \(\pi_A^{\rm hi}\) and \(\pi_{\max}\) do not depend on the slot. The prior \(\pi_A^{\rm lo}\) depends on it
through \(u\).

\begin{paperresult}{Corollary A.2 (Coexistence at every prior)}
Take the example, any slot with \(0<\lambda<1\) and \(0<\beta<1/2\), any prior \(\pi\in(0,1)\), and the default tie
rules. The set \(\mathcal C(\pi)\) of strengths on the path at which a quiet market and an informative market coexist
is the set of Proposition~3(a) on the path,
\[
\mathcal C(\pi)=\bigl\{r\in[r_0,r_2]:\ r\ge r_A,\ \ \pi<\tau(r)\le M(\pi),\ \ \rho\DT(r)\max\{\pi,1-\pi\}\le k\bigr\}.
\]
\begin{enumerate}[label=(\alph*),leftmargin=2em]
\item If \(\pi<\pi_A^{\rm lo}\), \(\mathcal C(\pi)\) is empty. At \(\pi=\pi_A^{\rm lo}\), \(\mathcal C(\pi)=\{r_A\}\).
\item If \(\pi_A^{\rm lo}<\pi<\pi_A^{\rm hi}\), then \(\mathcal C(\pi)=[r_A,\min\{r_N(\pi),r_C(\pi)\}]\), a nondegenerate
interval. Coexistence starts at \(r_A\) for every prior in this window and every slot.
\item If \(\pi_A^{\rm hi}\le\pi<\pi_{\max}\), then \(\mathcal C(\pi)=(r_\pi,\min\{r_N(\pi),r_C(\pi)\}]\), a nonempty
interval with \(r_\pi\ge r_A\). Coexistence starts where the prior reaches the threshold.
\item If \(\pi\ge\pi_{\max}\), \(\mathcal C(\pi)\) is empty.
\end{enumerate}
At each \(r\in\mathcal C(\pi)\) the informative market is a mixed purchase that recruits or a full purchase, and the
threshold pair at \(r_A\) in case (b). Its entry is \(\rho+(1-\rho)\min\{(1-\pi)\,b\DT(r)/k,\ b+\pi\lambda\}\), which
rises strictly while the purchase is mixed and is then constant.
\end{paperresult}

Why the set has this form: by Proposition~3(a), coexistence at \(r\) needs \(\pi<\tau(r)\le M(\pi)\) and
\(k/(1-\tau(r))\le\DT(r)\le\DT^{\rm N}\). Under \(\tau(r)\le M(\pi)<1\), the first spread condition is \(J(r)\ge k\), which
holds exactly when \(r\ge r_A\), because \(J\) rises strictly on the path. The number \(r_A\) contains neither the prior
nor the slot. The second spread condition is \(\rho\DT(r)\max\{\pi,1-\pi\}\le k\). Each case then follows from the
direction of each condition along the path.
\begin{itemize}
\item \emph{Low priors, case (a).} A low prior gives a low top posterior. Coexistence needs \(r\ge r_A\), where
\(\tau(r)\ge\tau_A\), and \(\tau(r)\le M(\pi)\). Both hold together only if \(M(\pi)\ge\tau_A\), that is
\(\pi\ge\pi_A^{\rm lo}\). In the paper's words, below \(\pi_A^{\rm lo}\) even the top posterior cannot reach the
threshold once a purchase can pay for itself.
\item \emph{Middle priors, case (b).} The prior lies below \(\tau_A\), so \(\pi<\tau(r)\) holds at every \(r\ge r_A\),
and the top posterior exceeds \(\tau_A\). Coexistence starts at \(r_A\), for every prior in the window, and ends when
either the threshold passes the top posterior or an order against a silent price starts to pay. At \(r_A\) the
informative market is the threshold pair.
\item \emph{Higher priors, case (c).} The prior is at least \(\tau_A\), so the threshold must first rise above the
prior. Coexistence starts just after \(r_\pi\), where the threshold passes the prior. The left end is open, because at
\(\tau=\pi\) the default challenger rule gives entry one after every order in the quiet market.
\item \emph{High priors, case (d).} For \(\pi\ge\pi_{\max}>1/2\), the quiet market needs \(\rho\DT(r)\pi\le k\). Any
\(r\) with \(\pi<\tau(r)\) has \(r>r_\pi\ge r_*\), so \(\rho\DT(r)\pi>\rho\DT(r_*)\tau(r_*)=k\). In the paper's words,
the quiet market ends before the prior falls below the threshold.
\end{itemize}
This agrees with the bound of part (a): coexistence needs \(\pi<1/(1+\rho)\), so a prior that strongly favours class
\(H\) rules it out. The paper reports the three priors at the example's slot in decimals; Chapter~\ref{app:numbers}
lists them.

\begin{chaptersummary}
Against one incumbent, a quiet market and an informative market can both be equilibria, because recruitment raises
the return of a purchase and informed purchases make recruitment happen. Under the default tie rules they coexist
exactly when \(\pi<\tau\le M(\pi)\) and \(k/(1-\tau)\le\DT\le k/(\rho\max\{\pi,1-\pi\})\): the threshold lies above
the prior and within reach of a purchase, a purchase that just reaches the threshold and recruits covers \(k\), and no
order against a silent price does. In the region the informative market is unique: class \(L\) abstains, only a
purchase recruits, and class \(H\) mixes until the doubt left after a purchase pays exactly \(k\). Entry there is
\(\rho+(1-\rho)\min\{(1-\pi)b\DT/k,\,b+\pi\lambda\}\), which rises strictly with the spread while class \(H\) mixes and
is constant once it buys for sure. Along a path the sure-purchase strengths form one interval, and the whole
coexistence set is one interval under the balance condition, which uniform incumbents with an affordable high class
satisfy. In the example at the prior \(1/2\), coexistence runs exactly from \(r_A\) to \(r_N\), with a sure purchase
from \(r_F\), in the order \(r_0<r_A<r_F<r_N<r_J<r_S<r_1<r_C<r_2\).
\end{chaptersummary}

% =====================================================================
\needspace{10\baselineskip}
\answersection

\answer{\ref{q:coexistence-binds}}{Section~\ref{sec:coexistence-quiet} gives \(K_H=\rho\DT(1-\pi)\) and
\(K_L=\rho\DT\pi\). Then
\begin{align*}
K_L-K_H &= \rho\,\DT\,\bigl[\pi-(1-\pi)\bigr] \why{subtract}\\
&= \rho\,\DT\,(2\pi-1)>0 \why{\(\pi>1/2\); \(\rho\DT>0\)}
\end{align*}
So \(K_L>K_H\), and \(\max\{\pi,1-\pi\}=\pi\): the sale of class \(L\) binds, and \(\DT^{\rm N}=k/(\rho\pi)\). Class
\(L\), which the prior disfavours, is closer to trade.

In words: the silent price keeps the belief at \(\pi\). Class \(L\) earns on the weight that the price gives to class
\(H\), which is \(\pi\), large. Class \(H\) earns on the weight on class \(L\), which is \(1-\pi\), small. The trading
bound is different. It is a floor over every candidate, so it puts each order's posterior as far toward the
investor's own class as the slot allows: \(M\) after a purchase, \(m\) after a sale. When the prior already leans
toward class \(H\), a purchase can push the belief close to one, and class \(H\) is left with little doubt to earn
on. In the quiet market no order moves the belief, so the lean of the prior helps class \(L\) instead.}

\answer{\ref{q:coexistence-mix}}{At \(\pi=1/2\), the formula of Section~\ref{sec:coexistence-informative} gives
\begin{align*}
\sigma_H(1) &= \frac{1}{\pi u}\Bigl[\frac{(1-\pi)\DT}{k}-1\Bigr] \why{the purchase probability while class \(H\) mixes}\\
&= \frac{2}{u}\Bigl[\frac{\DT}{2k}-1\Bigr] \why{\(\pi=1-\pi=1/2\)}\\
&= \frac{1}{u}\Bigl[\frac{\DT}{k}-2\Bigr] \why{multiply out the \(2\)}
\end{align*}
The slope in \(\DT\) is \(1/(uk)>0\), so \(\sigma_H(1)\) rises strictly with the spread. It equals one exactly when
\(\DT/k-2=u\), that is \(\DT=k(2+u)\). At \(\pi=1/2\), \(1-M(1/2)=1/(2+u)\)
(Section~\ref{sec:coexistence-example}), so
\[
\DT^{\rm F}=\frac{k}{1-M(1/2)}=k\,(2+u).
\]
So class \(H\) buys for sure exactly from the spread \(\DT^{\rm F}\) on, as Step~4 of
Section~\ref{sec:coexistence-informative} says.}

\answer{\ref{q:coexistence-edge}}{On the left edge \(J=\DT(1-\tau)=k\), so \(\DT=k/(1-\tau)\) and \(\DT(1-M)<\DT(1-\tau)=k\)
because \(\tau<M\). So \(\DT<\DT^{\rm F}\), and the rate is the second term of \eqref{eq:coexistence-rate}:
\begin{align*}
s_+ &= \frac{b}{\pi}\Bigl[\frac{(1-\pi)\DT}{k}-1\Bigr] \why{class \(H\) mixes}\\
&= \frac{b}{\pi}\Bigl[\frac{1-\pi}{1-\tau}-1\Bigr] \why{\(\DT/k=1/(1-\tau)\)}\\
&= \frac{b}{\pi}\cdot\frac{1-\pi-1+\tau}{1-\tau}=\frac{b\,(\tau-\pi)}{\pi\,(1-\tau)}=s^* \why{common denominator}
\end{align*}
At this rate the posterior after a purchase is exactly \(\tau\). The high-cost challenger is indifferent, and the
purchase earns \(e(1)\DT(1-\tau)=e(1)\,k\). Class \(H\) mixes, so this must equal \(k\), which needs \(e(1)=1\): every
high-cost challenger must prepare at the tie. The default challenger rule gives exactly this. Under the other
challenger rule, \(e(1)\) may be any value in \([\rho,1]\), and the pair (T) of Lemma~A.7 requires \(e(1)=k/J\). On the
edge \(J=k\), this is again \(e(1)=1\). So the market exists under both challenger rules here, but under the other rule
it is one of many entry values the rule allows, and only the value one keeps class \(H\) indifferent. Inside the region,
where \(J>k\), the other rule adds a threshold pair with the smaller entry \(k/J\).}

\answer{\ref{q:coexistence-gap}}{Write \(\delta>0\) for the common fall of \(g_H\) and \(g_L\) between two strengths \(r<r'\)
of the stretch. Then \(\gamma(r')=g_H(r)-\delta-[g_L(r)-\delta]=\gamma(r)\), and
\begin{align*}
\tau(r')-\tau(r) &= \frac{c_H-g_L(r)+\delta}{\gamma(r)}-\frac{c_H-g_L(r)}{\gamma(r)} \why{definition of \(\tau\); \(\gamma\) fixed}\\
&= \frac{\delta}{\gamma(r)}>0 \why{\(\delta>0\), \(\gamma>0\)}\\[3pt]
J(r')-J(r) &= \DT\,[1-\tau(r')]-\DT\,[1-\tau(r)] \why{\(\DT\) fixed}\\
&= -\DT\,\bigl[\tau(r')-\tau(r)\bigr]<0 \why{\(\DT>0\); the line above}
\end{align*}
So \(\tau\) rises and \(J\) falls. The product in the balance condition changes by
\[
\DT\,[g_H(r')-c_H]-\DT\,[g_H(r)-c_H]=-\DT\,\delta<0,
\]
so it falls, and the balance condition fails. The condition \(J\ge k\) can switch off on this stretch: if \(J(r)\ge k\)
at the start, \(J\) can fall below \(k\). Later, on a stretch where the spread widens, \(J\) can rise above \(k\) again.
The conditions \(\pi<\tau\), \(\tau\le M\) and \(\DT\le\DT^{\rm N}\) cannot switch back, because each moves in one
direction along every path (Section~\ref{sec:coexistence-proof-b}). This is the two-interval path of
Section~\ref{sec:coexistence-race}.}

\answer{\ref{q:coexistence-abstain}}{\emph{Step 1: the two inequalities.} On \((r_N,r_S]\), \(r>r_N\) gives
\(\rho\DT(r)/2>k\), and \(r\le r_S\) gives \(\rho\DT(r)m\le k\), because \(\DT\) rises strictly on the path. So
\(\rho\DT m\le k<\rho\DT\pi\) at \(\pi=1/2\).

\emph{Step 2: Lemma~A.8.} On the path \(\tau>1/2=\pi\), so the lemma applies, and every sale has entry \(\rho\). Its
last sentence covers this case: when an indifferent investor abstains, no rate qualifies. In detail:
\begin{itemize}
\item \(s_-=0\) fails: \(K_L=\rho\DT/2>k\), so a sale strictly beats abstention and class \(L\) must sell;
\item \(s_-=\lambda\) fails: \(K_L=\rho\DT m\le k\), and under this rule the investor sells for sure only if
\(K_L>k\);
\item an interior rate fails: it needs a mix, and this rule allows none.
\end{itemize}

\emph{Step 3: the full purchase does not help.} By Lemma~A.3(d), class \(L\)'s payoffs involve only \(s_-\),
\(\mu(-1)\) and \(e(-1)\). With no wrong-signed orders, \(\mu(-1)=\mu_-(s_-)\) depends only on \(s_-\), and
\(e(-1)=\rho\). So class \(H\)'s purchase rate does not enter class \(L\)'s problem, and no choice of class \(H\) can
give class \(L\) a best response. Every candidate fails (E5) for class \(L\), so the equilibrium set is empty on
\((r_N,r_S]\). In words: the low class wants to sell when nobody expects it to, and does not want to sell for sure when
everybody expects it to; only a mix would balance the two, and this rule forbids a mix.}
